\documentclass[11pt]{article}

    \usepackage[breakable]{tcolorbox}
    \usepackage{parskip} % Stop auto-indenting (to mimic markdown behaviour)
    

    % Basic figure setup, for now with no caption control since it's done
    % automatically by Pandoc (which extracts ![](path) syntax from Markdown).
    \usepackage{graphicx}
    % Maintain compatibility with old templates. Remove in nbconvert 6.0
    \let\Oldincludegraphics\includegraphics
    % Ensure that by default, figures have no caption (until we provide a
    % proper Figure object with a Caption API and a way to capture that
    % in the conversion process - todo).
    \usepackage{caption}
    \DeclareCaptionFormat{nocaption}{}
    \captionsetup{format=nocaption,aboveskip=0pt,belowskip=0pt}

    \usepackage{float}
    \floatplacement{figure}{H} % forces figures to be placed at the correct location
    \usepackage{xcolor} % Allow colors to be defined
    \usepackage{enumerate} % Needed for markdown enumerations to work
    \usepackage{geometry} % Used to adjust the document margins
    \usepackage{amsmath} % Equations
    \usepackage{amssymb} % Equations
    \usepackage{textcomp} % defines textquotesingle
    % Hack from http://tex.stackexchange.com/a/47451/13684:
    \AtBeginDocument{%
        \def\PYZsq{\textquotesingle}% Upright quotes in Pygmentized code
    }
    \usepackage{upquote} % Upright quotes for verbatim code
    \usepackage{eurosym} % defines \euro

    \usepackage{iftex}
    \ifPDFTeX
        \usepackage[T1]{fontenc}
        \IfFileExists{alphabeta.sty}{
              \usepackage{alphabeta}
          }{
              \usepackage[mathletters]{ucs}
              \usepackage[utf8x]{inputenc}
          }
    \else
        \usepackage{fontspec}
        \usepackage{unicode-math}
    \fi

    \usepackage{fancyvrb} % verbatim replacement that allows latex
    \usepackage{grffile} % extends the file name processing of package graphics
                         % to support a larger range
    \makeatletter % fix for old versions of grffile with XeLaTeX
    \@ifpackagelater{grffile}{2019/11/01}
    {
      % Do nothing on new versions
    }
    {
      \def\Gread@@xetex#1{%
        \IfFileExists{"\Gin@base".bb}%
        {\Gread@eps{\Gin@base.bb}}%
        {\Gread@@xetex@aux#1}%
      }
    }
    \makeatother
    \usepackage[Export]{adjustbox} % Used to constrain images to a maximum size
    \adjustboxset{max size={0.9\linewidth}{0.9\paperheight}}

    % The hyperref package gives us a pdf with properly built
    % internal navigation ('pdf bookmarks' for the table of contents,
    % internal cross-reference links, web links for URLs, etc.)
    \usepackage{hyperref}
    % The default LaTeX title has an obnoxious amount of whitespace. By default,
    % titling removes some of it. It also provides customization options.
    \usepackage{titling}
    \usepackage{longtable} % longtable support required by pandoc >1.10
    \usepackage{booktabs}  % table support for pandoc > 1.12.2
    \usepackage{array}     % table support for pandoc >= 2.11.3
    \usepackage{calc}      % table minipage width calculation for pandoc >= 2.11.1
    \usepackage[inline]{enumitem} % IRkernel/repr support (it uses the enumerate* environment)
    \usepackage[normalem]{ulem} % ulem is needed to support strikethroughs (\sout)
                                % normalem makes italics be italics, not underlines
    \usepackage{mathrsfs}
    

    
    % Colors for the hyperref package
    \definecolor{urlcolor}{rgb}{0,.145,.698}
    \definecolor{linkcolor}{rgb}{.71,0.21,0.01}
    \definecolor{citecolor}{rgb}{.12,.54,.11}

    % ANSI colors
    \definecolor{ansi-black}{HTML}{3E424D}
    \definecolor{ansi-black-intense}{HTML}{282C36}
    \definecolor{ansi-red}{HTML}{E75C58}
    \definecolor{ansi-red-intense}{HTML}{B22B31}
    \definecolor{ansi-green}{HTML}{00A250}
    \definecolor{ansi-green-intense}{HTML}{007427}
    \definecolor{ansi-yellow}{HTML}{DDB62B}
    \definecolor{ansi-yellow-intense}{HTML}{B27D12}
    \definecolor{ansi-blue}{HTML}{208FFB}
    \definecolor{ansi-blue-intense}{HTML}{0065CA}
    \definecolor{ansi-magenta}{HTML}{D160C4}
    \definecolor{ansi-magenta-intense}{HTML}{A03196}
    \definecolor{ansi-cyan}{HTML}{60C6C8}
    \definecolor{ansi-cyan-intense}{HTML}{258F8F}
    \definecolor{ansi-white}{HTML}{C5C1B4}
    \definecolor{ansi-white-intense}{HTML}{A1A6B2}
    \definecolor{ansi-default-inverse-fg}{HTML}{FFFFFF}
    \definecolor{ansi-default-inverse-bg}{HTML}{000000}

    % common color for the border for error outputs.
    \definecolor{outerrorbackground}{HTML}{FFDFDF}

    % commands and environments needed by pandoc snippets
    % extracted from the output of `pandoc -s`
    \providecommand{\tightlist}{%
      \setlength{\itemsep}{0pt}\setlength{\parskip}{0pt}}
    \DefineVerbatimEnvironment{Highlighting}{Verbatim}{commandchars=\\\{\}}
    % Add ',fontsize=\small' for more characters per line
    \newenvironment{Shaded}{}{}
    \newcommand{\KeywordTok}[1]{\textcolor[rgb]{0.00,0.44,0.13}{\textbf{{#1}}}}
    \newcommand{\DataTypeTok}[1]{\textcolor[rgb]{0.56,0.13,0.00}{{#1}}}
    \newcommand{\DecValTok}[1]{\textcolor[rgb]{0.25,0.63,0.44}{{#1}}}
    \newcommand{\BaseNTok}[1]{\textcolor[rgb]{0.25,0.63,0.44}{{#1}}}
    \newcommand{\FloatTok}[1]{\textcolor[rgb]{0.25,0.63,0.44}{{#1}}}
    \newcommand{\CharTok}[1]{\textcolor[rgb]{0.25,0.44,0.63}{{#1}}}
    \newcommand{\StringTok}[1]{\textcolor[rgb]{0.25,0.44,0.63}{{#1}}}
    \newcommand{\CommentTok}[1]{\textcolor[rgb]{0.38,0.63,0.69}{\textit{{#1}}}}
    \newcommand{\OtherTok}[1]{\textcolor[rgb]{0.00,0.44,0.13}{{#1}}}
    \newcommand{\AlertTok}[1]{\textcolor[rgb]{1.00,0.00,0.00}{\textbf{{#1}}}}
    \newcommand{\FunctionTok}[1]{\textcolor[rgb]{0.02,0.16,0.49}{{#1}}}
    \newcommand{\RegionMarkerTok}[1]{{#1}}
    \newcommand{\ErrorTok}[1]{\textcolor[rgb]{1.00,0.00,0.00}{\textbf{{#1}}}}
    \newcommand{\NormalTok}[1]{{#1}}

    % Additional commands for more recent versions of Pandoc
    \newcommand{\ConstantTok}[1]{\textcolor[rgb]{0.53,0.00,0.00}{{#1}}}
    \newcommand{\SpecialCharTok}[1]{\textcolor[rgb]{0.25,0.44,0.63}{{#1}}}
    \newcommand{\VerbatimStringTok}[1]{\textcolor[rgb]{0.25,0.44,0.63}{{#1}}}
    \newcommand{\SpecialStringTok}[1]{\textcolor[rgb]{0.73,0.40,0.53}{{#1}}}
    \newcommand{\ImportTok}[1]{{#1}}
    \newcommand{\DocumentationTok}[1]{\textcolor[rgb]{0.73,0.13,0.13}{\textit{{#1}}}}
    \newcommand{\AnnotationTok}[1]{\textcolor[rgb]{0.38,0.63,0.69}{\textbf{\textit{{#1}}}}}
    \newcommand{\CommentVarTok}[1]{\textcolor[rgb]{0.38,0.63,0.69}{\textbf{\textit{{#1}}}}}
    \newcommand{\VariableTok}[1]{\textcolor[rgb]{0.10,0.09,0.49}{{#1}}}
    \newcommand{\ControlFlowTok}[1]{\textcolor[rgb]{0.00,0.44,0.13}{\textbf{{#1}}}}
    \newcommand{\OperatorTok}[1]{\textcolor[rgb]{0.40,0.40,0.40}{{#1}}}
    \newcommand{\BuiltInTok}[1]{{#1}}
    \newcommand{\ExtensionTok}[1]{{#1}}
    \newcommand{\PreprocessorTok}[1]{\textcolor[rgb]{0.74,0.48,0.00}{{#1}}}
    \newcommand{\AttributeTok}[1]{\textcolor[rgb]{0.49,0.56,0.16}{{#1}}}
    \newcommand{\InformationTok}[1]{\textcolor[rgb]{0.38,0.63,0.69}{\textbf{\textit{{#1}}}}}
    \newcommand{\WarningTok}[1]{\textcolor[rgb]{0.38,0.63,0.69}{\textbf{\textit{{#1}}}}}


    % Define a nice break command that doesn't care if a line doesn't already
    % exist.
    \def\br{\hspace*{\fill} \\* }
    % Math Jax compatibility definitions
    \def\gt{>}
    \def\lt{<}
    \let\Oldtex\TeX
    \let\Oldlatex\LaTeX
    \renewcommand{\TeX}{\textrm{\Oldtex}}
    \renewcommand{\LaTeX}{\textrm{\Oldlatex}}
    % Document parameters
    % Document title
    \title{HW-5}
    
    
    
    
    
% Pygments definitions
\makeatletter
\def\PY@reset{\let\PY@it=\relax \let\PY@bf=\relax%
    \let\PY@ul=\relax \let\PY@tc=\relax%
    \let\PY@bc=\relax \let\PY@ff=\relax}
\def\PY@tok#1{\csname PY@tok@#1\endcsname}
\def\PY@toks#1+{\ifx\relax#1\empty\else%
    \PY@tok{#1}\expandafter\PY@toks\fi}
\def\PY@do#1{\PY@bc{\PY@tc{\PY@ul{%
    \PY@it{\PY@bf{\PY@ff{#1}}}}}}}
\def\PY#1#2{\PY@reset\PY@toks#1+\relax+\PY@do{#2}}

\@namedef{PY@tok@w}{\def\PY@tc##1{\textcolor[rgb]{0.73,0.73,0.73}{##1}}}
\@namedef{PY@tok@c}{\let\PY@it=\textit\def\PY@tc##1{\textcolor[rgb]{0.24,0.48,0.48}{##1}}}
\@namedef{PY@tok@cp}{\def\PY@tc##1{\textcolor[rgb]{0.61,0.40,0.00}{##1}}}
\@namedef{PY@tok@k}{\let\PY@bf=\textbf\def\PY@tc##1{\textcolor[rgb]{0.00,0.50,0.00}{##1}}}
\@namedef{PY@tok@kp}{\def\PY@tc##1{\textcolor[rgb]{0.00,0.50,0.00}{##1}}}
\@namedef{PY@tok@kt}{\def\PY@tc##1{\textcolor[rgb]{0.69,0.00,0.25}{##1}}}
\@namedef{PY@tok@o}{\def\PY@tc##1{\textcolor[rgb]{0.40,0.40,0.40}{##1}}}
\@namedef{PY@tok@ow}{\let\PY@bf=\textbf\def\PY@tc##1{\textcolor[rgb]{0.67,0.13,1.00}{##1}}}
\@namedef{PY@tok@nb}{\def\PY@tc##1{\textcolor[rgb]{0.00,0.50,0.00}{##1}}}
\@namedef{PY@tok@nf}{\def\PY@tc##1{\textcolor[rgb]{0.00,0.00,1.00}{##1}}}
\@namedef{PY@tok@nc}{\let\PY@bf=\textbf\def\PY@tc##1{\textcolor[rgb]{0.00,0.00,1.00}{##1}}}
\@namedef{PY@tok@nn}{\let\PY@bf=\textbf\def\PY@tc##1{\textcolor[rgb]{0.00,0.00,1.00}{##1}}}
\@namedef{PY@tok@ne}{\let\PY@bf=\textbf\def\PY@tc##1{\textcolor[rgb]{0.80,0.25,0.22}{##1}}}
\@namedef{PY@tok@nv}{\def\PY@tc##1{\textcolor[rgb]{0.10,0.09,0.49}{##1}}}
\@namedef{PY@tok@no}{\def\PY@tc##1{\textcolor[rgb]{0.53,0.00,0.00}{##1}}}
\@namedef{PY@tok@nl}{\def\PY@tc##1{\textcolor[rgb]{0.46,0.46,0.00}{##1}}}
\@namedef{PY@tok@ni}{\let\PY@bf=\textbf\def\PY@tc##1{\textcolor[rgb]{0.44,0.44,0.44}{##1}}}
\@namedef{PY@tok@na}{\def\PY@tc##1{\textcolor[rgb]{0.41,0.47,0.13}{##1}}}
\@namedef{PY@tok@nt}{\let\PY@bf=\textbf\def\PY@tc##1{\textcolor[rgb]{0.00,0.50,0.00}{##1}}}
\@namedef{PY@tok@nd}{\def\PY@tc##1{\textcolor[rgb]{0.67,0.13,1.00}{##1}}}
\@namedef{PY@tok@s}{\def\PY@tc##1{\textcolor[rgb]{0.73,0.13,0.13}{##1}}}
\@namedef{PY@tok@sd}{\let\PY@it=\textit\def\PY@tc##1{\textcolor[rgb]{0.73,0.13,0.13}{##1}}}
\@namedef{PY@tok@si}{\let\PY@bf=\textbf\def\PY@tc##1{\textcolor[rgb]{0.64,0.35,0.47}{##1}}}
\@namedef{PY@tok@se}{\let\PY@bf=\textbf\def\PY@tc##1{\textcolor[rgb]{0.67,0.36,0.12}{##1}}}
\@namedef{PY@tok@sr}{\def\PY@tc##1{\textcolor[rgb]{0.64,0.35,0.47}{##1}}}
\@namedef{PY@tok@ss}{\def\PY@tc##1{\textcolor[rgb]{0.10,0.09,0.49}{##1}}}
\@namedef{PY@tok@sx}{\def\PY@tc##1{\textcolor[rgb]{0.00,0.50,0.00}{##1}}}
\@namedef{PY@tok@m}{\def\PY@tc##1{\textcolor[rgb]{0.40,0.40,0.40}{##1}}}
\@namedef{PY@tok@gh}{\let\PY@bf=\textbf\def\PY@tc##1{\textcolor[rgb]{0.00,0.00,0.50}{##1}}}
\@namedef{PY@tok@gu}{\let\PY@bf=\textbf\def\PY@tc##1{\textcolor[rgb]{0.50,0.00,0.50}{##1}}}
\@namedef{PY@tok@gd}{\def\PY@tc##1{\textcolor[rgb]{0.63,0.00,0.00}{##1}}}
\@namedef{PY@tok@gi}{\def\PY@tc##1{\textcolor[rgb]{0.00,0.52,0.00}{##1}}}
\@namedef{PY@tok@gr}{\def\PY@tc##1{\textcolor[rgb]{0.89,0.00,0.00}{##1}}}
\@namedef{PY@tok@ge}{\let\PY@it=\textit}
\@namedef{PY@tok@gs}{\let\PY@bf=\textbf}
\@namedef{PY@tok@gp}{\let\PY@bf=\textbf\def\PY@tc##1{\textcolor[rgb]{0.00,0.00,0.50}{##1}}}
\@namedef{PY@tok@go}{\def\PY@tc##1{\textcolor[rgb]{0.44,0.44,0.44}{##1}}}
\@namedef{PY@tok@gt}{\def\PY@tc##1{\textcolor[rgb]{0.00,0.27,0.87}{##1}}}
\@namedef{PY@tok@err}{\def\PY@bc##1{{\setlength{\fboxsep}{\string -\fboxrule}\fcolorbox[rgb]{1.00,0.00,0.00}{1,1,1}{\strut ##1}}}}
\@namedef{PY@tok@kc}{\let\PY@bf=\textbf\def\PY@tc##1{\textcolor[rgb]{0.00,0.50,0.00}{##1}}}
\@namedef{PY@tok@kd}{\let\PY@bf=\textbf\def\PY@tc##1{\textcolor[rgb]{0.00,0.50,0.00}{##1}}}
\@namedef{PY@tok@kn}{\let\PY@bf=\textbf\def\PY@tc##1{\textcolor[rgb]{0.00,0.50,0.00}{##1}}}
\@namedef{PY@tok@kr}{\let\PY@bf=\textbf\def\PY@tc##1{\textcolor[rgb]{0.00,0.50,0.00}{##1}}}
\@namedef{PY@tok@bp}{\def\PY@tc##1{\textcolor[rgb]{0.00,0.50,0.00}{##1}}}
\@namedef{PY@tok@fm}{\def\PY@tc##1{\textcolor[rgb]{0.00,0.00,1.00}{##1}}}
\@namedef{PY@tok@vc}{\def\PY@tc##1{\textcolor[rgb]{0.10,0.09,0.49}{##1}}}
\@namedef{PY@tok@vg}{\def\PY@tc##1{\textcolor[rgb]{0.10,0.09,0.49}{##1}}}
\@namedef{PY@tok@vi}{\def\PY@tc##1{\textcolor[rgb]{0.10,0.09,0.49}{##1}}}
\@namedef{PY@tok@vm}{\def\PY@tc##1{\textcolor[rgb]{0.10,0.09,0.49}{##1}}}
\@namedef{PY@tok@sa}{\def\PY@tc##1{\textcolor[rgb]{0.73,0.13,0.13}{##1}}}
\@namedef{PY@tok@sb}{\def\PY@tc##1{\textcolor[rgb]{0.73,0.13,0.13}{##1}}}
\@namedef{PY@tok@sc}{\def\PY@tc##1{\textcolor[rgb]{0.73,0.13,0.13}{##1}}}
\@namedef{PY@tok@dl}{\def\PY@tc##1{\textcolor[rgb]{0.73,0.13,0.13}{##1}}}
\@namedef{PY@tok@s2}{\def\PY@tc##1{\textcolor[rgb]{0.73,0.13,0.13}{##1}}}
\@namedef{PY@tok@sh}{\def\PY@tc##1{\textcolor[rgb]{0.73,0.13,0.13}{##1}}}
\@namedef{PY@tok@s1}{\def\PY@tc##1{\textcolor[rgb]{0.73,0.13,0.13}{##1}}}
\@namedef{PY@tok@mb}{\def\PY@tc##1{\textcolor[rgb]{0.40,0.40,0.40}{##1}}}
\@namedef{PY@tok@mf}{\def\PY@tc##1{\textcolor[rgb]{0.40,0.40,0.40}{##1}}}
\@namedef{PY@tok@mh}{\def\PY@tc##1{\textcolor[rgb]{0.40,0.40,0.40}{##1}}}
\@namedef{PY@tok@mi}{\def\PY@tc##1{\textcolor[rgb]{0.40,0.40,0.40}{##1}}}
\@namedef{PY@tok@il}{\def\PY@tc##1{\textcolor[rgb]{0.40,0.40,0.40}{##1}}}
\@namedef{PY@tok@mo}{\def\PY@tc##1{\textcolor[rgb]{0.40,0.40,0.40}{##1}}}
\@namedef{PY@tok@ch}{\let\PY@it=\textit\def\PY@tc##1{\textcolor[rgb]{0.24,0.48,0.48}{##1}}}
\@namedef{PY@tok@cm}{\let\PY@it=\textit\def\PY@tc##1{\textcolor[rgb]{0.24,0.48,0.48}{##1}}}
\@namedef{PY@tok@cpf}{\let\PY@it=\textit\def\PY@tc##1{\textcolor[rgb]{0.24,0.48,0.48}{##1}}}
\@namedef{PY@tok@c1}{\let\PY@it=\textit\def\PY@tc##1{\textcolor[rgb]{0.24,0.48,0.48}{##1}}}
\@namedef{PY@tok@cs}{\let\PY@it=\textit\def\PY@tc##1{\textcolor[rgb]{0.24,0.48,0.48}{##1}}}

\def\PYZbs{\char`\\}
\def\PYZus{\char`\_}
\def\PYZob{\char`\{}
\def\PYZcb{\char`\}}
\def\PYZca{\char`\^}
\def\PYZam{\char`\&}
\def\PYZlt{\char`\<}
\def\PYZgt{\char`\>}
\def\PYZsh{\char`\#}
\def\PYZpc{\char`\%}
\def\PYZdl{\char`\$}
\def\PYZhy{\char`\-}
\def\PYZsq{\char`\'}
\def\PYZdq{\char`\"}
\def\PYZti{\char`\~}
% for compatibility with earlier versions
\def\PYZat{@}
\def\PYZlb{[}
\def\PYZrb{]}
\makeatother


    % For linebreaks inside Verbatim environment from package fancyvrb.
    \makeatletter
        \newbox\Wrappedcontinuationbox
        \newbox\Wrappedvisiblespacebox
        \newcommand*\Wrappedvisiblespace {\textcolor{red}{\textvisiblespace}}
        \newcommand*\Wrappedcontinuationsymbol {\textcolor{red}{\llap{\tiny$\m@th\hookrightarrow$}}}
        \newcommand*\Wrappedcontinuationindent {3ex }
        \newcommand*\Wrappedafterbreak {\kern\Wrappedcontinuationindent\copy\Wrappedcontinuationbox}
        % Take advantage of the already applied Pygments mark-up to insert
        % potential linebreaks for TeX processing.
        %        {, <, #, %, $, ' and ": go to next line.
        %        _, }, ^, &, >, - and ~: stay at end of broken line.
        % Use of \textquotesingle for straight quote.
        \newcommand*\Wrappedbreaksatspecials {%
            \def\PYGZus{\discretionary{\char`\_}{\Wrappedafterbreak}{\char`\_}}%
            \def\PYGZob{\discretionary{}{\Wrappedafterbreak\char`\{}{\char`\{}}%
            \def\PYGZcb{\discretionary{\char`\}}{\Wrappedafterbreak}{\char`\}}}%
            \def\PYGZca{\discretionary{\char`\^}{\Wrappedafterbreak}{\char`\^}}%
            \def\PYGZam{\discretionary{\char`\&}{\Wrappedafterbreak}{\char`\&}}%
            \def\PYGZlt{\discretionary{}{\Wrappedafterbreak\char`\<}{\char`\<}}%
            \def\PYGZgt{\discretionary{\char`\>}{\Wrappedafterbreak}{\char`\>}}%
            \def\PYGZsh{\discretionary{}{\Wrappedafterbreak\char`\#}{\char`\#}}%
            \def\PYGZpc{\discretionary{}{\Wrappedafterbreak\char`\%}{\char`\%}}%
            \def\PYGZdl{\discretionary{}{\Wrappedafterbreak\char`\$}{\char`\$}}%
            \def\PYGZhy{\discretionary{\char`\-}{\Wrappedafterbreak}{\char`\-}}%
            \def\PYGZsq{\discretionary{}{\Wrappedafterbreak\textquotesingle}{\textquotesingle}}%
            \def\PYGZdq{\discretionary{}{\Wrappedafterbreak\char`\"}{\char`\"}}%
            \def\PYGZti{\discretionary{\char`\~}{\Wrappedafterbreak}{\char`\~}}%
        }
        % Some characters . , ; ? ! / are not pygmentized.
        % This macro makes them "active" and they will insert potential linebreaks
        \newcommand*\Wrappedbreaksatpunct {%
            \lccode`\~`\.\lowercase{\def~}{\discretionary{\hbox{\char`\.}}{\Wrappedafterbreak}{\hbox{\char`\.}}}%
            \lccode`\~`\,\lowercase{\def~}{\discretionary{\hbox{\char`\,}}{\Wrappedafterbreak}{\hbox{\char`\,}}}%
            \lccode`\~`\;\lowercase{\def~}{\discretionary{\hbox{\char`\;}}{\Wrappedafterbreak}{\hbox{\char`\;}}}%
            \lccode`\~`\:\lowercase{\def~}{\discretionary{\hbox{\char`\:}}{\Wrappedafterbreak}{\hbox{\char`\:}}}%
            \lccode`\~`\?\lowercase{\def~}{\discretionary{\hbox{\char`\?}}{\Wrappedafterbreak}{\hbox{\char`\?}}}%
            \lccode`\~`\!\lowercase{\def~}{\discretionary{\hbox{\char`\!}}{\Wrappedafterbreak}{\hbox{\char`\!}}}%
            \lccode`\~`\/\lowercase{\def~}{\discretionary{\hbox{\char`\/}}{\Wrappedafterbreak}{\hbox{\char`\/}}}%
            \catcode`\.\active
            \catcode`\,\active
            \catcode`\;\active
            \catcode`\:\active
            \catcode`\?\active
            \catcode`\!\active
            \catcode`\/\active
            \lccode`\~`\~
        }
    \makeatother

    \let\OriginalVerbatim=\Verbatim
    \makeatletter
    \renewcommand{\Verbatim}[1][1]{%
        %\parskip\z@skip
        \sbox\Wrappedcontinuationbox {\Wrappedcontinuationsymbol}%
        \sbox\Wrappedvisiblespacebox {\FV@SetupFont\Wrappedvisiblespace}%
        \def\FancyVerbFormatLine ##1{\hsize\linewidth
            \vtop{\raggedright\hyphenpenalty\z@\exhyphenpenalty\z@
                \doublehyphendemerits\z@\finalhyphendemerits\z@
                \strut ##1\strut}%
        }%
        % If the linebreak is at a space, the latter will be displayed as visible
        % space at end of first line, and a continuation symbol starts next line.
        % Stretch/shrink are however usually zero for typewriter font.
        \def\FV@Space {%
            \nobreak\hskip\z@ plus\fontdimen3\font minus\fontdimen4\font
            \discretionary{\copy\Wrappedvisiblespacebox}{\Wrappedafterbreak}
            {\kern\fontdimen2\font}%
        }%

        % Allow breaks at special characters using \PYG... macros.
        \Wrappedbreaksatspecials
        % Breaks at punctuation characters . , ; ? ! and / need catcode=\active
        \OriginalVerbatim[#1,codes*=\Wrappedbreaksatpunct]%
    }
    \makeatother

    % Exact colors from NB
    \definecolor{incolor}{HTML}{303F9F}
    \definecolor{outcolor}{HTML}{D84315}
    \definecolor{cellborder}{HTML}{CFCFCF}
    \definecolor{cellbackground}{HTML}{F7F7F7}

    % prompt
    \makeatletter
    \newcommand{\boxspacing}{\kern\kvtcb@left@rule\kern\kvtcb@boxsep}
    \makeatother
    \newcommand{\prompt}[4]{
        {\ttfamily\llap{{\color{#2}[#3]:\hspace{3pt}#4}}\vspace{-\baselineskip}}
    }
    

    
    % Prevent overflowing lines due to hard-to-break entities
    \sloppy
    % Setup hyperref package
    \hypersetup{
      breaklinks=true,  % so long urls are correctly broken across lines
      colorlinks=true,
      urlcolor=urlcolor,
      linkcolor=linkcolor,
      citecolor=citecolor,
      }
    % Slightly bigger margins than the latex defaults
    
    \geometry{verbose,tmargin=1in,bmargin=1in,lmargin=1in,rmargin=1in}
    
    

\begin{document}
    
    \maketitle
    
    

    
    \hypertarget{homework-5}{%
\subsection{Homework 5}\label{homework-5}}

\textbf{Your name}.

    \begin{tcolorbox}[breakable, size=fbox, boxrule=1pt, pad at break*=1mm,colback=cellbackground, colframe=cellborder]
\prompt{In}{incolor}{2}{\boxspacing}
\begin{Verbatim}[commandchars=\\\{\}]
\PY{k}{using}\PY{+w}{ }\PY{n}{LinearAlgebra}\PY{p}{,}\PY{+w}{ }\PY{n}{Plots}
\PY{k}{using}\PY{+w}{ }\PY{n}{LaTeXStrings}
\end{Verbatim}
\end{tcolorbox}

    \hypertarget{q-7.2.5}{%
\subsection{Q 7.2.5}\label{q-7.2.5}}

Let \(A\) be a 2x2 matrix defined below.

\begin{enumerate}
\def\labelenumi{(\alph{enumi})}
\item
  Find the eigenvalue decomposition of \(A\) using \texttt{eigen}. The
  result is a data structure with two components, \texttt{.values} and
  \texttt{.vectors}
\item
  Write a function \texttt{p\_matrix} that uses Eq. (7.2.7) to evaluate
  \(p(A)\), where \(p(x) = (x - \pi)^4\). You can find a matrix inverse
  using the function \texttt{inv} and you can construct a matrix from
  its diagonal entries using \texttt{diagm}.
\item
  Write a function \texttt{cos\_matrix} that uses the function analog of
  (7.2.7) to evaluate \(p(A) = \cos(A)\).
\end{enumerate}

    \begin{tcolorbox}[breakable, size=fbox, boxrule=1pt, pad at break*=1mm,colback=cellbackground, colframe=cellborder]
\prompt{In}{incolor}{3}{\boxspacing}
\begin{Verbatim}[commandchars=\\\{\}]
\PY{n}{A}\PY{+w}{ }\PY{o}{=}\PY{+w}{ }\PY{p}{[}\PY{+w}{ }\PY{l+m+mi}{4}\PY{+w}{ }\PY{l+m+mi}{1}\PY{p}{;}\PY{+w}{ }\PY{l+m+mi}{4}\PY{+w}{ }\PY{l+m+mi}{4}\PY{p}{]}\PY{+w}{ }\PY{o}{*}\PY{+w}{ }\PY{n+nb}{pi}\PY{+w}{ }\PY{o}{/}\PY{+w}{ }\PY{l+m+mi}{6}
\end{Verbatim}
\end{tcolorbox}

    
    \begin{Verbatim}[commandchars=\\\{\}]
2×2 Matrix\{Float64\}:
 2.0944  0.523599
 2.0944  2.0944
    \end{Verbatim}

    
    Compare the numerical value of your solution to the ``long way'',
computing the fourth power of \(A - \pi\) explicitly. Also compare to
Julia's built-in algorithm for computing cosine of a matrix.

    \begin{tcolorbox}[breakable, size=fbox, boxrule=1pt, pad at break*=1mm,colback=cellbackground, colframe=cellborder]
\prompt{In}{incolor}{4}{\boxspacing}
\begin{Verbatim}[commandchars=\\\{\}]
\PY{p}{(}\PY{n}{A}\PY{+w}{ }\PY{o}{\PYZhy{}}\PY{+w}{ }\PY{n}{one}\PY{p}{(}\PY{n}{A}\PY{p}{)}\PY{+w}{ }\PY{o}{*}\PY{+w}{ }\PY{n+nb}{pi}\PY{p}{)}\PY{+w}{ }\PY{o}{\PYZca{}}\PY{+w}{ }\PY{l+m+mi}{4}\PY{+w}{ }\PY{c}{\PYZsh{} built\PYZhy{}in calculation \PYZhy{}\PYZhy{} the notation hides the algorithm}
\end{Verbatim}
\end{tcolorbox}

    
    \begin{Verbatim}[commandchars=\\\{\}]
2×2 Matrix\{Float64\}:
   9.62065  -4.81033
 -19.2413    9.62065
    \end{Verbatim}

    
    \begin{tcolorbox}[breakable, size=fbox, boxrule=1pt, pad at break*=1mm,colback=cellbackground, colframe=cellborder]
\prompt{In}{incolor}{5}{\boxspacing}
\begin{Verbatim}[commandchars=\\\{\}]
\PY{n}{cos}\PY{p}{(}\PY{n}{A}\PY{p}{)}\PY{+w}{ }\PY{c}{\PYZsh{} built\PYZhy{}in}
\end{Verbatim}
\end{tcolorbox}

    
    \begin{Verbatim}[commandchars=\\\{\}]
2×2 Matrix\{Float64\}:
 -0.25  -0.375
 -1.5   -0.25
    \end{Verbatim}

    
    \begin{tcolorbox}[breakable, size=fbox, boxrule=1pt, pad at break*=1mm,colback=cellbackground, colframe=cellborder]
\prompt{In}{incolor}{6}{\boxspacing}
\begin{Verbatim}[commandchars=\\\{\}]
\PY{n}{ev}\PY{+w}{ }\PY{o}{=}\PY{+w}{ }\PY{n}{eigen}\PY{p}{(}\PY{n}{A}\PY{p}{)}
\end{Verbatim}
\end{tcolorbox}

    
    \begin{Verbatim}[commandchars=\\\{\}]
Eigen\{Float64, Float64, Matrix\{Float64\}, Vector\{Float64\}\}
values:
2-element Vector\{Float64\}:
 1.0471975511965976
 3.141592653589793
vectors:
2×2 Matrix\{Float64\}:
 -0.447214  0.447214
  0.894427  0.894427
    \end{Verbatim}

    
    \emph{Solution}

    \begin{enumerate}
\def\labelenumi{(\alph{enumi})}
\tightlist
\item
  The eigenvalue decomposition is
\end{enumerate}

\[
    A = P\Lambda P^{-1},\quad 
    \Lambda = \begin{bmatrix}
        \lambda_1&0\\0&\lambda_2
    \end{bmatrix} = \begin{bmatrix}
        1.0472&0\\0&3.1416
    \end{bmatrix},\quad P = \begin{bmatrix}
        \mathbf{v_1}&\mathbf{v_2}
    \end{bmatrix} = \begin{bmatrix}
        -0.4472&0.4472\\0.8944&0.8944
    \end{bmatrix}
\]

and
\(p^{-1}=\begin{bmatrix}  -1.1180&0.5590\\1.1180&0.5590  \end{bmatrix}\),
thus

\[
    A = \begin{bmatrix}
        -0.4472&0.4472\\0.8944&0.8944
    \end{bmatrix}\begin{bmatrix}
        1.0472&0\\0&3.1416
    \end{bmatrix}\begin{bmatrix}
        -1.1180&0.5590\\1.1180&0.5590
    \end{bmatrix}
\]

    \begin{tcolorbox}[breakable, size=fbox, boxrule=1pt, pad at break*=1mm,colback=cellbackground, colframe=cellborder]
\prompt{In}{incolor}{7}{\boxspacing}
\begin{Verbatim}[commandchars=\\\{\}]
\PY{c}{\PYZsh{} (a)}
\PY{n}{ev}\PY{+w}{ }\PY{o}{=}\PY{+w}{ }\PY{n}{eigen}\PY{p}{(}\PY{n}{A}\PY{p}{)}
\PY{n}{display}\PY{p}{(}\PY{n}{ev}\PY{p}{)}
\PY{n}{display}\PY{p}{(}\PY{n}{inv}\PY{p}{(}\PY{n}{ev}\PY{o}{.}\PY{n}{vectors}\PY{p}{)}\PY{p}{)}
\PY{n}{display}\PY{p}{(}\PY{n}{ev}\PY{o}{.}\PY{n}{vectors}\PY{+w}{ }\PY{o}{*}\PY{+w}{ }\PY{n}{diagm}\PY{p}{(}\PY{n}{ev}\PY{o}{.}\PY{n}{values}\PY{p}{)}\PY{+w}{ }\PY{o}{*}\PY{+w}{ }\PY{n}{inv}\PY{p}{(}\PY{n}{ev}\PY{o}{.}\PY{n}{vectors}\PY{p}{)}\PY{p}{)}

\PY{c}{\PYZsh{} verify}
\PY{n}{ev}\PY{o}{.}\PY{n}{vectors}\PY{+w}{ }\PY{o}{*}\PY{+w}{ }\PY{n}{diagm}\PY{p}{(}\PY{n}{ev}\PY{o}{.}\PY{n}{values}\PY{p}{)}\PY{+w}{ }\PY{o}{*}\PY{+w}{ }\PY{n}{inv}\PY{p}{(}\PY{n}{ev}\PY{o}{.}\PY{n}{vectors}\PY{p}{)}\PY{+w}{ }\PY{o}{\PYZhy{}}\PY{+w}{ }\PY{n}{A}
\end{Verbatim}
\end{tcolorbox}

    
    \begin{Verbatim}[commandchars=\\\{\}]
Eigen\{Float64, Float64, Matrix\{Float64\}, Vector\{Float64\}\}
values:
2-element Vector\{Float64\}:
 1.0471975511965976
 3.141592653589793
vectors:
2×2 Matrix\{Float64\}:
 -0.447214  0.447214
  0.894427  0.894427
    \end{Verbatim}

    
    
    \begin{Verbatim}[commandchars=\\\{\}]
2×2 Matrix\{Float64\}:
 -1.11803  0.559017
  1.11803  0.559017
    \end{Verbatim}

    
    
    \begin{Verbatim}[commandchars=\\\{\}]
2×2 Matrix\{Float64\}:
 2.0944  0.523599
 2.0944  2.0944
    \end{Verbatim}

    
    
    \begin{Verbatim}[commandchars=\\\{\}]
2×2 Matrix\{Float64\}:
 0.0  1.11022e-16
 0.0  4.44089e-16
    \end{Verbatim}

    
    \begin{enumerate}
\def\labelenumi{(\alph{enumi})}
\setcounter{enumi}{1}
\tightlist
\item
  \(P^{-1}(A-\pi)^4P=(P^{-1}(A-\pi)P)^4=(\Lambda-\pi)^4\), hence
  \((A-\pi)^4=P(\Lambda-\pi)^4P^{-1}\).
\end{enumerate}

    \begin{tcolorbox}[breakable, size=fbox, boxrule=1pt, pad at break*=1mm,colback=cellbackground, colframe=cellborder]
\prompt{In}{incolor}{8}{\boxspacing}
\begin{Verbatim}[commandchars=\\\{\}]
\PY{p}{(}\PY{n}{ev}\PY{o}{.}\PY{n}{vectors}\PY{p}{)}\PY{+w}{ }\PY{o}{*}\PY{+w}{ }\PY{n}{diagm}\PY{p}{(}\PY{p}{(}\PY{n}{ev}\PY{o}{.}\PY{n}{values}\PY{+w}{ }\PY{o}{.\PYZhy{}}\PY{+w}{ }\PY{n+nb}{pi}\PY{p}{)}\PY{+w}{ }\PY{o}{.\PYZca{}}\PY{+w}{ }\PY{l+m+mi}{4}\PY{p}{)}\PY{+w}{ }\PY{o}{*}\PY{+w}{ }\PY{n}{inv}\PY{p}{(}\PY{n}{ev}\PY{o}{.}\PY{n}{vectors}\PY{p}{)}
\end{Verbatim}
\end{tcolorbox}

    
    \begin{Verbatim}[commandchars=\\\{\}]
2×2 Matrix\{Float64\}:
   9.62065  -4.81033
 -19.2413    9.62065
    \end{Verbatim}

    
    \begin{enumerate}
\def\labelenumi{(\alph{enumi})}
\setcounter{enumi}{2}
\tightlist
\item
  Let
  \(p(x)=\cos x=1-\dfrac{x^2}{2!}+\dfrac{x^4}{4!}-\dfrac{x^6}{6!}+\cdots\),
  then
\end{enumerate}

\[
    \begin{align}
        p(A)=\cos A &= I-\dfrac{A^2}{2!}+\dfrac{A^4}{4!}-\dfrac{A^6}{6!}+\cdots\\
        & = V\text{diag}(p(\lambda_1), p(\lambda_2), \cdots, p(\lambda_n))V^{-1}\\
        & = V\text{diag}(\cos(\lambda_1), \cos(\lambda_2), \cdots, \cos(\lambda_n))V^{-1}
    \end{align}
\]

    \begin{tcolorbox}[breakable, size=fbox, boxrule=1pt, pad at break*=1mm,colback=cellbackground, colframe=cellborder]
\prompt{In}{incolor}{9}{\boxspacing}
\begin{Verbatim}[commandchars=\\\{\}]
\PY{n}{ev}\PY{o}{.}\PY{n}{vectors}\PY{+w}{ }\PY{o}{*}\PY{+w}{ }\PY{n}{diagm}\PY{p}{(}\PY{n}{cos}\PY{o}{.}\PY{p}{(}\PY{n}{ev}\PY{o}{.}\PY{n}{values}\PY{p}{)}\PY{p}{)}\PY{+w}{ }\PY{o}{*}\PY{+w}{ }\PY{n}{inv}\PY{p}{(}\PY{n}{ev}\PY{o}{.}\PY{n}{vectors}\PY{p}{)}
\end{Verbatim}
\end{tcolorbox}

    
    \begin{Verbatim}[commandchars=\\\{\}]
2×2 Matrix\{Float64\}:
 -0.25  -0.375
 -1.5   -0.25
    \end{Verbatim}

    
    \hypertarget{q-7.2.8}{%
\subsection{Q 7.2.8}\label{q-7.2.8}}

Eigenvalues of random matrices and their perturbations can be very
interesting.

\begin{enumerate}
\def\labelenumi{(\alph{enumi})}
\item
  Let A=randn(60,60). Scatter plot its eigenvalues in the complex plane,
  using aspect\_ratio=1 and red diamonds as markers.
\item
  Let \(E\) be another random \(60\times 60\) matrix, and on top of the
  previous graph, plot the eigenvalues of \(A+0.05E\) as blue dots.
  Repeat this for 100 different values of \(E\).
\item
  Let T=triu(A). On a new graph, scatter plot the eigenvalues of \(T\)
  in the complex plane. (They all lie on the real axis.)
\item
  Repeat part (b) with \(T\) in place of \(A\).
\end{enumerate}

    \begin{tcolorbox}[breakable, size=fbox, boxrule=1pt, pad at break*=1mm,colback=cellbackground, colframe=cellborder]
\prompt{In}{incolor}{10}{\boxspacing}
\begin{Verbatim}[commandchars=\\\{\}]
\PY{c}{\PYZsh{} (a)}

\PY{n}{A}\PY{+w}{ }\PY{o}{=}\PY{+w}{ }\PY{n}{randn}\PY{p}{(}\PY{l+m+mi}{60}\PY{p}{,}\PY{+w}{ }\PY{l+m+mi}{60}\PY{p}{)}\PY{p}{;}
\PY{n}{ev}\PY{+w}{ }\PY{o}{=}\PY{+w}{ }\PY{n}{eigen}\PY{p}{(}\PY{n}{A}\PY{p}{)}
\PY{n}{scatter}\PY{p}{(}\PY{n}{ev}\PY{o}{.}\PY{n}{values}\PY{p}{,}\PY{+w}{ }\PY{n}{shape}\PY{+w}{ }\PY{o}{=}\PY{+w}{ }\PY{l+s+ss}{:diamond}\PY{p}{,}\PY{+w}{ }\PY{n}{color}\PY{+w}{ }\PY{o}{=}\PY{+w}{ }\PY{l+s+ss}{:red}\PY{p}{,}\PY{+w}{ }\PY{n}{label}\PY{+w}{ }\PY{o}{=}\PY{+w}{ }\PY{l+s}{\PYZdq{}}\PY{l+s}{\PYZdq{}}\PY{p}{,}\PY{+w}{ }\PY{n}{aspect\PYZus{}ratio}\PY{+w}{ }\PY{o}{=}\PY{+w}{ }\PY{l+m+mi}{1}\PY{p}{)}
\end{Verbatim}
\end{tcolorbox}

    \begin{center}
    \adjustimage{max size={0.9\linewidth}{0.9\paperheight}}{HW-5_files/HW-5_16_0.pdf}
    \end{center}
    { \hspace*{\fill} \\}
    
    \begin{tcolorbox}[breakable, size=fbox, boxrule=1pt, pad at break*=1mm,colback=cellbackground, colframe=cellborder]
\prompt{In}{incolor}{11}{\boxspacing}
\begin{Verbatim}[commandchars=\\\{\}]
\PY{c}{\PYZsh{} (b)}

\PY{n}{E}\PY{+w}{ }\PY{o}{=}\PY{+w}{ }\PY{n}{randn}\PY{p}{(}\PY{l+m+mi}{60}\PY{p}{,}\PY{+w}{ }\PY{l+m+mi}{60}\PY{p}{)}\PY{p}{;}
\PY{n}{ev2}\PY{+w}{ }\PY{o}{=}\PY{+w}{ }\PY{n}{eigen}\PY{p}{(}\PY{n}{A}\PY{+w}{ }\PY{o}{+}\PY{+w}{ }\PY{l+m+mf}{0.05}\PY{+w}{ }\PY{o}{*}\PY{+w}{ }\PY{n}{E}\PY{p}{)}
\PY{n}{display}\PY{p}{(}\PY{n}{scatter!}\PY{p}{(}\PY{n}{ev2}\PY{o}{.}\PY{n}{values}\PY{p}{,}\PY{+w}{ }\PY{n}{shape}\PY{+w}{ }\PY{o}{=}\PY{+w}{ }\PY{l+s+ss}{:circle}\PY{p}{,}\PY{+w}{ }\PY{n}{color}\PY{+w}{ }\PY{o}{=}\PY{+w}{ }\PY{l+s+ss}{:blue}\PY{p}{,}\PY{+w}{ }\PY{n}{label}\PY{+w}{ }\PY{o}{=}\PY{+w}{ }\PY{l+s}{\PYZdq{}}\PY{l+s}{\PYZdq{}}\PY{p}{)}\PY{p}{)}
\end{Verbatim}
\end{tcolorbox}

    \begin{center}
    \adjustimage{max size={0.9\linewidth}{0.9\paperheight}}{HW-5_files/HW-5_17_0.pdf}
    \end{center}
    { \hspace*{\fill} \\}
    
    \begin{tcolorbox}[breakable, size=fbox, boxrule=1pt, pad at break*=1mm,colback=cellbackground, colframe=cellborder]
\prompt{In}{incolor}{12}{\boxspacing}
\begin{Verbatim}[commandchars=\\\{\}]
\PY{c}{\PYZsh{} (c)}

\PY{n}{T}\PY{+w}{ }\PY{o}{=}\PY{+w}{ }\PY{n}{triu}\PY{p}{(}\PY{n}{A}\PY{p}{)}
\PY{n}{ev3}\PY{+w}{ }\PY{o}{=}\PY{+w}{ }\PY{n}{eigen}\PY{p}{(}\PY{n}{T}\PY{p}{)}
\PY{n}{scatter}\PY{p}{(}\PY{n}{real}\PY{p}{(}\PY{n}{ev3}\PY{o}{.}\PY{n}{values}\PY{p}{)}\PY{p}{,}\PY{+w}{ }\PY{n}{imag}\PY{p}{(}\PY{n}{ev3}\PY{o}{.}\PY{n}{values}\PY{p}{)}\PY{p}{,}\PY{+w}{ }\PY{n}{shape}\PY{+w}{ }\PY{o}{=}\PY{+w}{ }\PY{l+s+ss}{:diamond}\PY{p}{,}\PY{+w}{ }\PY{n}{color}\PY{+w}{ }\PY{o}{=}\PY{+w}{ }\PY{l+s+ss}{:orange}\PY{p}{,}\PY{+w}{ }\PY{n}{opacity}\PY{+w}{ }\PY{o}{=}\PY{+w}{ }\PY{l+m+mf}{0.4}\PY{p}{,}\PY{+w}{ }\PY{n}{label}\PY{+w}{ }\PY{o}{=}\PY{+w}{ }\PY{l+s}{\PYZdq{}}\PY{l+s}{\PYZdq{}}\PY{p}{,}\PY{+w}{ }\PY{n}{aspect\PYZus{}ratio}\PY{+w}{ }\PY{o}{=}\PY{+w}{ }\PY{l+m+mi}{1}\PY{p}{)}
\end{Verbatim}
\end{tcolorbox}

    \begin{center}
    \adjustimage{max size={0.9\linewidth}{0.9\paperheight}}{HW-5_files/HW-5_18_0.pdf}
    \end{center}
    { \hspace*{\fill} \\}
    
    \begin{tcolorbox}[breakable, size=fbox, boxrule=1pt, pad at break*=1mm,colback=cellbackground, colframe=cellborder]
\prompt{In}{incolor}{13}{\boxspacing}
\begin{Verbatim}[commandchars=\\\{\}]
\PY{c}{\PYZsh{} (d)}

\PY{n}{ev4}\PY{+w}{ }\PY{o}{=}\PY{+w}{ }\PY{n}{eigen}\PY{p}{(}\PY{n}{T}\PY{+w}{ }\PY{o}{+}\PY{+w}{ }\PY{l+m+mf}{0.05}\PY{+w}{ }\PY{o}{*}\PY{+w}{ }\PY{n}{E}\PY{p}{)}
\PY{n}{scatter!}\PY{p}{(}\PY{n}{real}\PY{p}{(}\PY{n}{ev3}\PY{o}{.}\PY{n}{values}\PY{p}{)}\PY{p}{,}\PY{+w}{ }\PY{n}{imag}\PY{p}{(}\PY{n}{ev3}\PY{o}{.}\PY{n}{values}\PY{p}{)}\PY{p}{,}\PY{+w}{ }\PY{n}{shape}\PY{+w}{ }\PY{o}{=}\PY{+w}{ }\PY{l+s+ss}{:circle}\PY{p}{,}\PY{+w}{ }\PY{n}{color}\PY{+w}{ }\PY{o}{=}\PY{+w}{ }\PY{l+s+ss}{:blue}\PY{p}{,}\PY{+w}{ }\PY{n}{opacity}\PY{+w}{ }\PY{o}{=}\PY{+w}{ }\PY{l+m+mf}{0.4}\PY{p}{,}\PY{+w}{ }\PY{n}{label}\PY{+w}{ }\PY{o}{=}\PY{+w}{ }\PY{l+s}{\PYZdq{}}\PY{l+s}{\PYZdq{}}\PY{p}{)}
\PY{n}{current}\PY{p}{(}\PY{p}{)}
\end{Verbatim}
\end{tcolorbox}

    \begin{center}
    \adjustimage{max size={0.9\linewidth}{0.9\paperheight}}{HW-5_files/HW-5_19_0.pdf}
    \end{center}
    { \hspace*{\fill} \\}
    
    \hypertarget{q-7.3.4}{%
\subsection{Q 7.3.4}\label{q-7.3.4}}

See Demo 7.1.7 for how to get the ``mandrill'' test image. Make a
log-linear scatter plot of the singular values of the matrix of
grayscale intensity values. (The shape of this graph is surprisingly
similar across a wide range of images.)

    \begin{tcolorbox}[breakable, size=fbox, boxrule=1pt, pad at break*=1mm,colback=cellbackground, colframe=cellborder]
\prompt{In}{incolor}{14}{\boxspacing}
\begin{Verbatim}[commandchars=\\\{\}]
\PY{k}{using}\PY{+w}{ }\PY{n}{Images}\PY{p}{,}\PY{+w}{ }\PY{n}{TestImages}
\end{Verbatim}
\end{tcolorbox}

    \emph{Solution}:

    \begin{tcolorbox}[breakable, size=fbox, boxrule=1pt, pad at break*=1mm,colback=cellbackground, colframe=cellborder]
\prompt{In}{incolor}{15}{\boxspacing}
\begin{Verbatim}[commandchars=\\\{\}]
\PY{n}{img}\PY{+w}{ }\PY{o}{=}\PY{+w}{ }\PY{n}{Gray}\PY{o}{.}\PY{p}{(}\PY{n}{testimage}\PY{p}{(}\PY{l+s}{\PYZdq{}}\PY{l+s}{mandrill}\PY{l+s}{\PYZdq{}}\PY{p}{)}\PY{p}{)}
\end{Verbatim}
\end{tcolorbox}

    \begin{center}
    \adjustimage{max size={0.9\linewidth}{0.9\paperheight}}{HW-5_files/HW-5_23_0.png}
    \end{center}
    { \hspace*{\fill} \\}
    
    \begin{tcolorbox}[breakable, size=fbox, boxrule=1pt, pad at break*=1mm,colback=cellbackground, colframe=cellborder]
\prompt{In}{incolor}{16}{\boxspacing}
\begin{Verbatim}[commandchars=\\\{\}]
\PY{n}{A}\PY{+w}{ }\PY{o}{=}\PY{+w}{ }\PY{k+kt}{Float64}\PY{o}{.}\PY{p}{(}\PY{n}{Gray}\PY{o}{.}\PY{p}{(}\PY{n}{img}\PY{p}{)}\PY{p}{)}
\end{Verbatim}
\end{tcolorbox}

    
    \begin{Verbatim}[commandchars=\\\{\}]
512×512 Matrix\{Float64\}:
 0.568627   0.219608   0.192157   0.34902   …  0.443137  0.6        0.694118
 0.454902   0.396078   0.156863   0.262745     0.509804  0.584314   0.490196
 0.301961   0.447059   0.180392   0.180392     0.392157  0.345098   0.317647
 0.278431   0.533333   0.372549   0.188235     0.286275  0.294118   0.27451
 0.341176   0.478431   0.517647   0.196078     0.231373  0.313725   0.337255
 0.321569   0.298039   0.623529   0.239216  …  0.333333  0.596078   0.713725
 0.168627   0.2        0.52549    0.509804     0.72549   0.776471   0.643137
 0.129412   0.2        0.435294   0.6          0.631373  0.470588   0.541176
 0.223529   0.188235   0.219608   0.458824     0.529412  0.611765   0.478431
 0.14902    0.12549    0.188235   0.25098      0.254902  0.317647   0.533333
 ⋮                                          ⋱            ⋮          
 0.466667   0.501961   0.564706   0.556863     0.341176  0.321569   0.333333
 0.490196   0.490196   0.52549    0.568627     0.341176  0.301961   0.243137
 0.443137   0.470588   0.513725   0.537255  …  0.364706  0.32549    0.337255
 0.411765   0.435294   0.517647   0.568627     0.321569  0.32549    0.32549
 0.376471   0.447059   0.541176   0.596078     0.360784  0.360784   0.333333
 0.435294   0.52549    0.596078   0.541176     0.388235  0.352941   0.34902
 0.552941   0.568627   0.552941   0.572549     0.32549   0.34902    0.305882
 0.611765   0.552941   0.513725   0.466667  …  0.309804  0.313725   0.270588
 0.0431373  0.0431373  0.0509804  0.054902     0.027451  0.0156863  0.0156863
    \end{Verbatim}

    
    \begin{tcolorbox}[breakable, size=fbox, boxrule=1pt, pad at break*=1mm,colback=cellbackground, colframe=cellborder]
\prompt{In}{incolor}{17}{\boxspacing}
\begin{Verbatim}[commandchars=\\\{\}]
\PY{n}{svals}\PY{+w}{ }\PY{o}{=}\PY{+w}{ }\PY{n}{svdvals}\PY{p}{(}\PY{n}{A}\PY{p}{)}
\PY{n}{scatter}\PY{p}{(}\PY{n}{svals}\PY{p}{,}\PY{n}{xaxis}\PY{o}{=}\PY{p}{(}\PY{l+s+sa}{L}\PY{l+s}{\PYZdq{}}\PY{l+s}{i}\PY{l+s}{\PYZdq{}}\PY{p}{)}\PY{p}{,}\PY{+w}{ }\PY{n}{yaxis}\PY{o}{=}\PY{p}{(}\PY{l+s+ss}{:log10}\PY{p}{,}\PY{l+s+sa}{L}\PY{l+s}{\PYZdq{}}\PY{l+s}{\PYZbs{}}\PY{l+s}{sigma\PYZus{}i}\PY{l+s}{\PYZdq{}}\PY{p}{)}\PY{p}{,}
\PY{n}{title}\PY{o}{=}\PY{l+s}{\PYZdq{}}\PY{l+s}{Singular values}\PY{l+s}{\PYZdq{}}\PY{p}{,}\PY{+w}{ }\PY{n}{label}\PY{o}{=}\PY{l+s}{\PYZdq{}}\PY{l+s}{\PYZdq{}}\PY{p}{)}
\end{Verbatim}
\end{tcolorbox}

    \begin{center}
    \adjustimage{max size={0.9\linewidth}{0.9\paperheight}}{HW-5_files/HW-5_25_0.pdf}
    \end{center}
    { \hspace*{\fill} \\}
    
    \hypertarget{q-7.4.4}{%
\subsection{Q 7.4.4}\label{q-7.4.4}}

The range of the function \(R_A(x)\) is a subset of the complex plane
known as the field of values of the matrix \(A\). Use 500 random
(possibly complex - you decide what is needed) vectors to plot points in
the field of values of \(A\). Then compute its eigenvalues and guess
what the exact field of values is.

    \begin{tcolorbox}[breakable, size=fbox, boxrule=1pt, pad at break*=1mm,colback=cellbackground, colframe=cellborder]
\prompt{In}{incolor}{18}{\boxspacing}
\begin{Verbatim}[commandchars=\\\{\}]
\PY{n}{R}\PY{p}{(}\PY{n}{A}\PY{p}{,}\PY{n}{x}\PY{p}{)}\PY{+w}{ }\PY{o}{=}\PY{+w}{ }\PY{n}{x}\PY{o}{\PYZsq{}}\PY{o}{*}\PY{n}{A}\PY{o}{*}\PY{n}{x}\PY{+w}{ }\PY{o}{/}\PY{+w}{ }\PY{p}{(}\PY{n}{x}\PY{o}{\PYZsq{}}\PY{+w}{ }\PY{o}{*}\PY{+w}{ }\PY{n}{x}\PY{p}{)}
\PY{n}{A}\PY{+w}{ }\PY{o}{=}\PY{+w}{ }\PY{p}{[}\PY{+w}{ }\PY{l+m+mi}{1}\PY{+w}{ }\PY{l+m+mi}{0}\PY{+w}{ }\PY{o}{\PYZhy{}}\PY{l+m+mi}{2}\PY{+w}{ }\PY{p}{;}\PY{+w}{ }\PY{l+m+mi}{0}\PY{+w}{ }\PY{l+m+mi}{2}\PY{+w}{ }\PY{l+m+mi}{0}\PY{+w}{ }\PY{p}{;}\PY{+w}{ }\PY{o}{\PYZhy{}}\PY{l+m+mi}{2}\PY{+w}{ }\PY{l+m+mi}{0}\PY{+w}{ }\PY{l+m+mi}{1}\PY{+w}{ }\PY{p}{]}
\end{Verbatim}
\end{tcolorbox}

    
    \begin{Verbatim}[commandchars=\\\{\}]
3×3 Matrix\{Int64\}:
  1  0  -2
  0  2   0
 -2  0   1
    \end{Verbatim}

    
    \emph{Solution}:

    \begin{tcolorbox}[breakable, size=fbox, boxrule=1pt, pad at break*=1mm,colback=cellbackground, colframe=cellborder]
\prompt{In}{incolor}{19}{\boxspacing}
\begin{Verbatim}[commandchars=\\\{\}]
\PY{n}{vals}\PY{+w}{ }\PY{o}{=}\PY{+w}{ }\PY{k+kt}{ComplexF64}\PY{p}{[}\PY{p}{]}
\PY{k}{for}\PY{+w}{ }\PY{n}{i}\PY{+w}{ }\PY{o}{=}\PY{+w}{ }\PY{l+m+mi}{1}\PY{o}{:}\PY{l+m+mi}{500}
\PY{+w}{    }\PY{n}{push!}\PY{p}{(}\PY{n}{vals}\PY{p}{,}\PY{+w}{ }\PY{n}{R}\PY{p}{(}\PY{n}{A}\PY{p}{,}\PY{+w}{ }\PY{n}{randn}\PY{p}{(}\PY{k+kt}{ComplexF64}\PY{p}{,}\PY{+w}{ }\PY{l+m+mi}{3}\PY{p}{)}\PY{p}{)}\PY{p}{)}
\PY{k}{end}

\PY{n}{scatter}\PY{p}{(}\PY{n}{vals}\PY{p}{,}\PY{+w}{ }\PY{n}{label}\PY{+w}{ }\PY{o}{=}\PY{+w}{ }\PY{l+s}{\PYZdq{}}\PY{l+s}{\PYZdq{}}\PY{p}{,}\PY{+w}{ }\PY{n}{mark}\PY{+w}{ }\PY{o}{=}\PY{+w}{ }\PY{l+s+ss}{:x}\PY{p}{,}\PY{+w}{ }\PY{n}{ylim}\PY{o}{=}\PY{p}{[}\PY{o}{\PYZhy{}}\PY{l+m+mf}{1e\PYZhy{}15}\PY{p}{,}\PY{l+m+mf}{1e\PYZhy{}15}\PY{p}{]}\PY{p}{)}
\end{Verbatim}
\end{tcolorbox}

    \begin{center}
    \adjustimage{max size={0.9\linewidth}{0.9\paperheight}}{HW-5_files/HW-5_29_0.pdf}
    \end{center}
    { \hspace*{\fill} \\}
    
    \begin{tcolorbox}[breakable, size=fbox, boxrule=1pt, pad at break*=1mm,colback=cellbackground, colframe=cellborder]
\prompt{In}{incolor}{20}{\boxspacing}
\begin{Verbatim}[commandchars=\\\{\}]
\PY{n}{eigen}\PY{p}{(}\PY{n}{A}\PY{p}{)}\PY{o}{.}\PY{n}{values}
\end{Verbatim}
\end{tcolorbox}

    
    \begin{Verbatim}[commandchars=\\\{\}]
3-element Vector\{Float64\}:
 -0.9999999999999978
  2.0
  3.0
    \end{Verbatim}

    
    The exact field of \(A\) should probably be \([-1,3]\).

    \hypertarget{q-7.5.5}{%
\subsection{Q 7.5.5}\label{q-7.5.5}}

Following Demo 7.5.2 as a guide, load the ``mandrill'' test image and
convert it to a matrix of floating-point pixel grayscale intensities.
Using the SVD, display as images the best approximations of rank 5, 10,
15, and 20.

    \emph{Solution}:

    \begin{tcolorbox}[breakable, size=fbox, boxrule=1pt, pad at break*=1mm,colback=cellbackground, colframe=cellborder]
\prompt{In}{incolor}{21}{\boxspacing}
\begin{Verbatim}[commandchars=\\\{\}]
\PY{n}{A}\PY{+w}{ }\PY{o}{=}\PY{+w}{ }\PY{k+kt}{Float64}\PY{o}{.}\PY{p}{(}\PY{n}{Gray}\PY{o}{.}\PY{p}{(}\PY{n}{testimage}\PY{p}{(}\PY{l+s}{\PYZdq{}}\PY{l+s}{mandrill}\PY{l+s}{\PYZdq{}}\PY{p}{)}\PY{p}{)}\PY{p}{)}
\end{Verbatim}
\end{tcolorbox}

    
    \begin{Verbatim}[commandchars=\\\{\}]
512×512 Matrix\{Float64\}:
 0.568627   0.219608   0.192157   0.34902   …  0.443137  0.6        0.694118
 0.454902   0.396078   0.156863   0.262745     0.509804  0.584314   0.490196
 0.301961   0.447059   0.180392   0.180392     0.392157  0.345098   0.317647
 0.278431   0.533333   0.372549   0.188235     0.286275  0.294118   0.27451
 0.341176   0.478431   0.517647   0.196078     0.231373  0.313725   0.337255
 0.321569   0.298039   0.623529   0.239216  …  0.333333  0.596078   0.713725
 0.168627   0.2        0.52549    0.509804     0.72549   0.776471   0.643137
 0.129412   0.2        0.435294   0.6          0.631373  0.470588   0.541176
 0.223529   0.188235   0.219608   0.458824     0.529412  0.611765   0.478431
 0.14902    0.12549    0.188235   0.25098      0.254902  0.317647   0.533333
 ⋮                                          ⋱            ⋮          
 0.466667   0.501961   0.564706   0.556863     0.341176  0.321569   0.333333
 0.490196   0.490196   0.52549    0.568627     0.341176  0.301961   0.243137
 0.443137   0.470588   0.513725   0.537255  …  0.364706  0.32549    0.337255
 0.411765   0.435294   0.517647   0.568627     0.321569  0.32549    0.32549
 0.376471   0.447059   0.541176   0.596078     0.360784  0.360784   0.333333
 0.435294   0.52549    0.596078   0.541176     0.388235  0.352941   0.34902
 0.552941   0.568627   0.552941   0.572549     0.32549   0.34902    0.305882
 0.611765   0.552941   0.513725   0.466667  …  0.309804  0.313725   0.270588
 0.0431373  0.0431373  0.0509804  0.054902     0.027451  0.0156863  0.0156863
    \end{Verbatim}

    
    \begin{tcolorbox}[breakable, size=fbox, boxrule=1pt, pad at break*=1mm,colback=cellbackground, colframe=cellborder]
\prompt{In}{incolor}{22}{\boxspacing}
\begin{Verbatim}[commandchars=\\\{\}]
\PY{n}{U}\PY{p}{,}\PY{n}{σ}\PY{p}{,}\PY{n}{V}\PY{+w}{ }\PY{o}{=}\PY{+w}{ }\PY{n}{svd}\PY{p}{(}\PY{n}{A}\PY{p}{)}
\PY{n}{scatter}\PY{p}{(}\PY{n}{σ}\PY{p}{,}\PY{n}{xaxis}\PY{o}{=}\PY{p}{(}\PY{l+s+sa}{L}\PY{l+s}{\PYZdq{}}\PY{l+s}{i}\PY{l+s}{\PYZdq{}}\PY{p}{)}\PY{p}{,}\PY{+w}{ }\PY{n}{yaxis}\PY{o}{=}\PY{p}{(}\PY{l+s+ss}{:log10}\PY{p}{,}\PY{l+s+sa}{L}\PY{l+s}{\PYZdq{}}\PY{l+s}{\PYZbs{}}\PY{l+s}{sigma\PYZus{}i}\PY{l+s}{\PYZdq{}}\PY{p}{)}\PY{p}{,}
\PY{+w}{    }\PY{n}{title}\PY{o}{=}\PY{l+s}{\PYZdq{}}\PY{l+s}{Singular values}\PY{l+s}{\PYZdq{}}\PY{p}{,}\PY{+w}{ }\PY{n}{label}\PY{o}{=}\PY{l+s}{\PYZdq{}}\PY{l+s}{\PYZdq{}}\PY{p}{)}
\end{Verbatim}
\end{tcolorbox}

    \begin{center}
    \adjustimage{max size={0.9\linewidth}{0.9\paperheight}}{HW-5_files/HW-5_35_0.pdf}
    \end{center}
    { \hspace*{\fill} \\}
    
    \begin{tcolorbox}[breakable, size=fbox, boxrule=1pt, pad at break*=1mm,colback=cellbackground, colframe=cellborder]
\prompt{In}{incolor}{23}{\boxspacing}
\begin{Verbatim}[commandchars=\\\{\}]
\PY{n}{plt}\PY{+w}{ }\PY{o}{=}\PY{+w}{ }\PY{n}{plot}\PY{p}{(}\PY{n}{layout}\PY{o}{=}\PY{p}{(}\PY{l+m+mi}{2}\PY{p}{,}\PY{l+m+mi}{2}\PY{p}{)}\PY{p}{,}\PY{n}{frame}\PY{o}{=}\PY{l+s+ss}{:none}\PY{p}{,}\PY{n}{aspect\PYZus{}ratio}\PY{o}{=}\PY{l+m+mi}{1}\PY{p}{,}\PY{n}{titlefontsize}\PY{o}{=}\PY{l+m+mi}{10}\PY{p}{)}
\PY{k}{for}\PY{+w}{ }\PY{n}{i}\PY{+w}{ }\PY{k}{in}\PY{+w}{ }\PY{l+m+mi}{1}\PY{o}{:}\PY{l+m+mi}{4}
\PY{+w}{    }\PY{n}{k}\PY{+w}{ }\PY{o}{=}\PY{+w}{ }\PY{l+m+mi}{5}\PY{n}{i}
\PY{+w}{    }\PY{n}{Ak}\PY{+w}{ }\PY{o}{=}\PY{+w}{ }\PY{n}{U}\PY{p}{[}\PY{o}{:}\PY{p}{,}\PY{l+m+mi}{1}\PY{o}{:}\PY{n}{k}\PY{p}{]}\PY{o}{*}\PY{n}{diagm}\PY{p}{(}\PY{n}{σ}\PY{p}{[}\PY{l+m+mi}{1}\PY{o}{:}\PY{n}{k}\PY{p}{]}\PY{p}{)}\PY{o}{*}\PY{n}{V}\PY{p}{[}\PY{o}{:}\PY{p}{,}\PY{l+m+mi}{1}\PY{o}{:}\PY{n}{k}\PY{p}{]}\PY{o}{\PYZsq{}}
\PY{+w}{    }\PY{n}{plot!}\PY{p}{(}\PY{n}{Gray}\PY{o}{.}\PY{p}{(}\PY{n}{Ak}\PY{p}{)}\PY{p}{,}\PY{n}{subplot}\PY{o}{=}\PY{n}{i}\PY{p}{,}\PY{n}{title}\PY{o}{=}\PY{l+s}{\PYZdq{}}\PY{l+s}{rank = }\PY{l+s+si}{\PYZdl{}k}\PY{l+s}{\PYZdq{}}\PY{p}{)}
\PY{k}{end}
\PY{n}{plt}
\end{Verbatim}
\end{tcolorbox}

    \begin{center}
    \adjustimage{max size={0.9\linewidth}{0.9\paperheight}}{HW-5_files/HW-5_36_0.pdf}
    \end{center}
    { \hspace*{\fill} \\}
    
    \hypertarget{q-9.5.2}{%
\subsection{Q 9.5.2}\label{q-9.5.2}}

\begin{enumerate}
\def\labelenumi{(\alph{enumi})}
\item
  Show that the functions \(\sin(r\pi x)\) and \(\sin(s\pi x)\) are
  identical at all of the nodes given in (9.5.2) if for an integer
  \(r-s=mN\). This important fact is called aliasing, and it implies
  that only finitely many frequencies can be distinguished on a fixed
  node set.
\item
  Demonstrate part (a) with a graph for the case \(N=11\), \(s=2\),
  \(r=-9\). Specifically, plot the two functions on one graph, and plot
  points to show that they intersect at all of the interpolation nodes.
\end{enumerate}

    \begin{tcolorbox}[breakable, size=fbox, boxrule=1pt, pad at break*=1mm,colback=cellbackground, colframe=cellborder]
\prompt{In}{incolor}{24}{\boxspacing}
\begin{Verbatim}[commandchars=\\\{\}]
\PY{n}{n}\PY{+w}{ }\PY{o}{=}\PY{+w}{ }\PY{l+m+mi}{5}
\PY{n}{N}\PY{+w}{ }\PY{o}{=}\PY{+w}{ }\PY{l+m+mi}{2}\PY{o}{*}\PY{n}{n}\PY{+w}{ }\PY{o}{+}\PY{+w}{ }\PY{l+m+mi}{1}
\PY{n}{tk}\PY{+w}{ }\PY{o}{=}\PY{+w}{ }\PY{l+m+mi}{2}\PY{o}{/}\PY{n}{N}\PY{+w}{ }\PY{o}{.*}\PY{+w}{ }\PY{p}{(}\PY{o}{\PYZhy{}}\PY{n}{n}\PY{o}{:}\PY{l+m+mi}{1}\PY{o}{:}\PY{n}{n}\PY{p}{)}\PY{p}{;}
\end{Verbatim}
\end{tcolorbox}

    \emph{Solution}:

    \begin{enumerate}
\def\labelenumi{(\alph{enumi})}
\tightlist
\item
  at \(x=x_k=\dfrac{2k}{N}\), and \(r-s=mN\), we have
\end{enumerate}

\[
    \sin(r\pi x)=\sin(\dfrac{2k(s+mN)\pi}{N})=\sin(\dfrac{2ks\pi}{N}+2km\pi)=\sin(\dfrac{2ks\pi}{N})=\sin(s\pi x)
\]

    \begin{tcolorbox}[breakable, size=fbox, boxrule=1pt, pad at break*=1mm,colback=cellbackground, colframe=cellborder]
\prompt{In}{incolor}{25}{\boxspacing}
\begin{Verbatim}[commandchars=\\\{\}]
\PY{c}{\PYZsh{} (b)}

\PY{n}{s}\PY{+w}{ }\PY{o}{=}\PY{+w}{ }\PY{l+m+mi}{2}\PY{p}{;}
\PY{n}{r}\PY{+w}{ }\PY{o}{=}\PY{+w}{ }\PY{o}{\PYZhy{}}\PY{l+m+mi}{9}\PY{p}{;}
\PY{n}{N}\PY{+w}{ }\PY{o}{=}\PY{+w}{ }\PY{l+m+mi}{11}\PY{p}{;}

\PY{n}{f1}\PY{p}{(}\PY{n}{x}\PY{p}{)}\PY{+w}{ }\PY{o}{=}\PY{+w}{ }\PY{n}{sin}\PY{p}{(}\PY{n}{r}\PY{+w}{ }\PY{o}{*}\PY{+w}{ }\PY{n+nb}{pi}\PY{+w}{ }\PY{o}{*}\PY{+w}{ }\PY{n}{x}\PY{p}{)}\PY{p}{;}
\PY{n}{f2}\PY{p}{(}\PY{n}{x}\PY{p}{)}\PY{+w}{ }\PY{o}{=}\PY{+w}{ }\PY{n}{sin}\PY{p}{(}\PY{n}{s}\PY{+w}{ }\PY{o}{*}\PY{+w}{ }\PY{n+nb}{pi}\PY{+w}{ }\PY{o}{*}\PY{+w}{ }\PY{n}{x}\PY{p}{)}\PY{p}{;}

\PY{n}{plot}\PY{p}{(}\PY{p}{[}\PY{n}{f1}\PY{p}{,}\PY{+w}{ }\PY{n}{f2}\PY{p}{]}\PY{p}{,}\PY{+w}{ }\PY{n}{xlim}\PY{+w}{ }\PY{o}{=}\PY{+w}{ }\PY{p}{[}\PY{n}{minimum}\PY{p}{(}\PY{n}{tk}\PY{p}{)}\PY{+w}{ }\PY{o}{\PYZhy{}}\PY{+w}{ }\PY{l+m+mf}{.1}\PY{p}{,}\PY{+w}{ }\PY{n}{maximum}\PY{p}{(}\PY{n}{tk}\PY{p}{)}\PY{+w}{ }\PY{o}{+}\PY{+w}{ }\PY{l+m+mf}{.1}\PY{p}{]}\PY{p}{,}\PY{+w}{ }\PY{n}{labels}\PY{+w}{ }\PY{o}{=}\PY{+w}{ }\PY{p}{[}\PY{l+s+sa}{L}\PY{l+s}{\PYZdq{}}\PY{l+s}{\PYZbs{}}\PY{l+s}{sin(r}\PY{l+s}{\PYZbs{}}\PY{l+s}{pi x)}\PY{l+s}{\PYZdq{}}\PY{+w}{ }\PY{l+s+sa}{L}\PY{l+s}{\PYZdq{}}\PY{l+s}{\PYZbs{}}\PY{l+s}{sin(s}\PY{l+s}{\PYZbs{}}\PY{l+s}{pi x)}\PY{l+s}{\PYZdq{}}\PY{p}{]}\PY{p}{)}
\PY{n}{scatter!}\PY{p}{(}\PY{n}{tk}\PY{p}{,}\PY{+w}{ }\PY{n}{f1}\PY{o}{.}\PY{p}{(}\PY{n}{tk}\PY{p}{)}\PY{p}{,}\PY{+w}{ }\PY{n}{label}\PY{+w}{ }\PY{o}{=}\PY{+w}{ }\PY{l+s+sa}{L}\PY{l+s}{\PYZdq{}}\PY{l+s}{t\PYZus{}k}\PY{l+s}{\PYZdq{}}\PY{p}{)}
\end{Verbatim}
\end{tcolorbox}

    \begin{center}
    \adjustimage{max size={0.9\linewidth}{0.9\paperheight}}{HW-5_files/HW-5_41_0.pdf}
    \end{center}
    { \hspace*{\fill} \\}
    
    \begin{tcolorbox}[breakable, size=fbox, boxrule=1pt, pad at break*=1mm,colback=cellbackground, colframe=cellborder]
\prompt{In}{incolor}{26}{\boxspacing}
\begin{Verbatim}[commandchars=\\\{\}]
\PY{n}{df}\PY{+w}{ }\PY{o}{=}\PY{+w}{ }\PY{p}{(}\PY{n}{x}\PY{p}{)}\PY{+w}{ }\PY{o}{\PYZhy{}\PYZgt{}}\PY{+w}{ }\PY{n}{f1}\PY{p}{(}\PY{n}{x}\PY{p}{)}\PY{+w}{ }\PY{o}{\PYZhy{}}\PY{+w}{ }\PY{n}{f2}\PY{p}{(}\PY{n}{x}\PY{p}{)}
\PY{n}{plot}\PY{p}{(}\PY{n}{df}\PY{p}{,}\PY{+w}{ }\PY{n}{xlim}\PY{+w}{ }\PY{o}{=}\PY{+w}{ }\PY{p}{[}\PY{n}{minimum}\PY{p}{(}\PY{n}{tk}\PY{p}{)}\PY{+w}{ }\PY{o}{\PYZhy{}}\PY{+w}{ }\PY{l+m+mf}{.1}\PY{p}{,}\PY{+w}{ }\PY{n}{maximum}\PY{p}{(}\PY{n}{tk}\PY{p}{)}\PY{+w}{ }\PY{o}{+}\PY{+w}{ }\PY{l+m+mf}{.1}\PY{p}{]}\PY{p}{,}\PY{+w}{ }\PY{n}{label}\PY{+w}{ }\PY{o}{=}\PY{+w}{ }\PY{l+s+sa}{L}\PY{l+s}{\PYZdq{}}\PY{l+s}{\PYZbs{}}\PY{l+s}{sin(r}\PY{l+s}{\PYZbs{}}\PY{l+s}{pi x) \PYZhy{} }\PY{l+s}{\PYZbs{}}\PY{l+s}{sin(s}\PY{l+s}{\PYZbs{}}\PY{l+s}{pi x)}\PY{l+s}{\PYZdq{}}\PY{p}{)}
\PY{n}{scatter!}\PY{p}{(}\PY{n}{tk}\PY{p}{,}\PY{+w}{ }\PY{n}{df}\PY{o}{.}\PY{p}{(}\PY{n}{tk}\PY{p}{)}\PY{p}{,}\PY{+w}{ }\PY{n}{label}\PY{+w}{ }\PY{o}{=}\PY{+w}{ }\PY{l+s+sa}{L}\PY{l+s}{\PYZdq{}}\PY{l+s}{t\PYZus{}k}\PY{l+s}{\PYZdq{}}\PY{p}{)}
\end{Verbatim}
\end{tcolorbox}

    \begin{center}
    \adjustimage{max size={0.9\linewidth}{0.9\paperheight}}{HW-5_files/HW-5_42_0.pdf}
    \end{center}
    { \hspace*{\fill} \\}
    
    \hypertarget{q-9.5.5}{%
\subsection{Q 9.5.5}\label{q-9.5.5}}

Spectral convergence is predicated on having infinitely many continuous
derivatives. At the other extreme is a function with a jump
discontinuity. Trigonometric interpolation across a jump leads to a lack
of convergence altogether, a fact famously known as the Gibbs
phenomenon.

\begin{enumerate}
\def\labelenumi{(\alph{enumi})}
\item
  Define \texttt{f(x)\ =\ sign(x+eps())}. This function jumps from \(0\)
  to \(1\) at \(-\epsilon\). Plot the function over \([-0.05, 0.15]\).
\item
  Let \(n=30\) and \(N=2n+1\). Using Function 9.5.3, add a plot of the
  trigonometric interpolant to \(f\) to the graph from part (a).
\item
  Repeat part (b) for \(n=80\) and \(n=180\).
\item
  You should see that the interpolants overshoot and oscillate near the
  step. The widths of the overshoots decrease with but the heights
  approach a limiting value. By zooming in to the graph, find the height
  of the overshoot to two decimal places.
\end{enumerate}

    \begin{tcolorbox}[breakable, size=fbox, boxrule=1pt, pad at break*=1mm,colback=cellbackground, colframe=cellborder]
\prompt{In}{incolor}{27}{\boxspacing}
\begin{Verbatim}[commandchars=\\\{\}]
\PY{c}{\PYZsh{} Function 9.5.3}
\PY{l+s}{\PYZdq{}\PYZdq{}\PYZdq{}}
\PY{l+s}{    triginterp(t,y)}

\PY{l+s}{Construct the trigonometric interpolant for the points defined by }
\PY{l+s}{vectors `t` and `y`.}
\PY{l+s}{\PYZdq{}\PYZdq{}\PYZdq{}}
\PY{k}{function}\PY{+w}{ }\PY{n}{triginterp}\PY{p}{(}\PY{n}{t}\PY{p}{,}\PY{n}{y}\PY{p}{)}
\PY{+w}{    }\PY{n}{N}\PY{+w}{ }\PY{o}{=}\PY{+w}{ }\PY{n}{length}\PY{p}{(}\PY{n}{t}\PY{p}{)}

\PY{+w}{    }\PY{k}{function}\PY{+w}{ }\PY{n}{τ}\PY{p}{(}\PY{n}{x}\PY{p}{)}
\PY{+w}{        }\PY{k}{if}\PY{+w}{ }\PY{n}{x}\PY{o}{==}\PY{l+m+mi}{0}
\PY{+w}{            }\PY{k}{return}\PY{+w}{ }\PY{l+m+mf}{1.0}
\PY{+w}{        }\PY{k}{else}
\PY{+w}{            }\PY{n}{denom}\PY{+w}{ }\PY{o}{=}\PY{+w}{ }\PY{n}{isodd}\PY{p}{(}\PY{n}{N}\PY{p}{)}\PY{+w}{ }\PY{o}{?}\PY{+w}{ }\PY{n}{N}\PY{o}{*}\PY{n}{sin}\PY{p}{(}\PY{n+nb}{π}\PY{o}{*}\PY{n}{x}\PY{o}{/}\PY{l+m+mi}{2}\PY{p}{)}\PY{+w}{ }\PY{o}{:}\PY{+w}{ }\PY{n}{N}\PY{o}{*}\PY{n}{tan}\PY{p}{(}\PY{n+nb}{π}\PY{o}{*}\PY{n}{x}\PY{o}{/}\PY{l+m+mi}{2}\PY{p}{)}
\PY{+w}{            }\PY{k}{return}\PY{+w}{ }\PY{n}{sin}\PY{p}{(}\PY{n}{N}\PY{o}{*}\PY{n+nb}{π}\PY{o}{*}\PY{n}{x}\PY{o}{/}\PY{l+m+mi}{2}\PY{p}{)}\PY{o}{/}\PY{n}{denom}
\PY{+w}{        }\PY{k}{end}
\PY{+w}{    }\PY{k}{end}

\PY{+w}{    }\PY{k}{return}\PY{+w}{ }\PY{k}{function}\PY{+w}{ }\PY{p}{(}\PY{n}{x}\PY{p}{)}
\PY{+w}{        }\PY{n}{sum}\PY{p}{(}\PY{+w}{ }\PY{n}{y}\PY{p}{[}\PY{n}{k}\PY{p}{]}\PY{o}{*}\PY{n}{τ}\PY{p}{(}\PY{n}{x}\PY{o}{\PYZhy{}}\PY{n}{t}\PY{p}{[}\PY{n}{k}\PY{p}{]}\PY{p}{)}\PY{+w}{ }\PY{k}{for}\PY{+w}{ }\PY{n}{k}\PY{+w}{ }\PY{k}{in}\PY{+w}{ }\PY{n}{eachindex}\PY{p}{(}\PY{n}{y}\PY{p}{)}\PY{+w}{ }\PY{p}{)}
\PY{+w}{    }\PY{k}{end}
\PY{k}{end}
\end{Verbatim}
\end{tcolorbox}

    
    \begin{Verbatim}[commandchars=\\\{\}]
triginterp
    \end{Verbatim}

    
    \emph{Solution}:

    \begin{tcolorbox}[breakable, size=fbox, boxrule=1pt, pad at break*=1mm,colback=cellbackground, colframe=cellborder]
\prompt{In}{incolor}{28}{\boxspacing}
\begin{Verbatim}[commandchars=\\\{\}]
\PY{n}{f}\PY{p}{(}\PY{n}{x}\PY{p}{)}\PY{+w}{ }\PY{o}{=}\PY{+w}{ }\PY{n}{sign}\PY{p}{(}\PY{n}{x}\PY{o}{+}\PY{n}{eps}\PY{p}{(}\PY{p}{)}\PY{p}{)}

\PY{n}{plot}\PY{p}{(}\PY{n}{f}\PY{p}{,}\PY{+w}{ }\PY{n}{xlim}\PY{+w}{ }\PY{o}{=}\PY{+w}{ }\PY{p}{[}\PY{o}{\PYZhy{}}\PY{l+m+mf}{0.05}\PY{p}{,}\PY{+w}{ }\PY{l+m+mf}{0.15}\PY{p}{]}\PY{p}{,}\PY{+w}{ }\PY{n}{label}\PY{+w}{ }\PY{o}{=}\PY{+w}{ }\PY{l+s+sa}{L}\PY{l+s}{\PYZdq{}}\PY{l+s}{\PYZbs{}}\PY{l+s}{mathrm\PYZob{}sign\PYZcb{}(x+}\PY{l+s+se}{\PYZbs{}v}\PY{l+s}{arepsilon)}\PY{l+s}{\PYZdq{}}\PY{p}{)}
\end{Verbatim}
\end{tcolorbox}

    \begin{center}
    \adjustimage{max size={0.9\linewidth}{0.9\paperheight}}{HW-5_files/HW-5_46_0.pdf}
    \end{center}
    { \hspace*{\fill} \\}
    
    \begin{tcolorbox}[breakable, size=fbox, boxrule=1pt, pad at break*=1mm,colback=cellbackground, colframe=cellborder]
\prompt{In}{incolor}{29}{\boxspacing}
\begin{Verbatim}[commandchars=\\\{\}]
\PY{n}{n}\PY{+w}{ }\PY{o}{=}\PY{+w}{ }\PY{l+m+mi}{30}\PY{p}{;}\PY{+w}{ }
\PY{n}{N}\PY{+w}{ }\PY{o}{=}\PY{+w}{ }\PY{l+m+mi}{2}\PY{+w}{ }\PY{o}{*}\PY{+w}{ }\PY{n}{n}\PY{+w}{ }\PY{o}{+}\PY{+w}{ }\PY{l+m+mi}{1}\PY{p}{;}
\PY{n}{ts}\PY{+w}{ }\PY{o}{=}\PY{+w}{ }\PY{l+m+mi}{2}\PY{o}{*}\PY{p}{(}\PY{o}{\PYZhy{}}\PY{n}{n}\PY{o}{:}\PY{n}{n}\PY{p}{)}\PY{o}{/}\PY{n}{N}
\PY{n}{ys}\PY{+w}{ }\PY{o}{=}\PY{+w}{ }\PY{n}{f}\PY{o}{.}\PY{p}{(}\PY{n}{ts}\PY{p}{)}
\PY{n}{fun}\PY{+w}{ }\PY{o}{=}\PY{+w}{ }\PY{n}{triginterp}\PY{p}{(}\PY{n}{ts}\PY{p}{,}\PY{+w}{ }\PY{n}{ys}\PY{p}{)}
\PY{n}{plot!}\PY{p}{(}\PY{n}{fun}\PY{p}{,}\PY{+w}{ }\PY{o}{\PYZhy{}}\PY{l+m+mi}{1}\PY{p}{,}\PY{+w}{ }\PY{l+m+mi}{1}\PY{p}{,}\PY{+w}{ }\PY{n}{label}\PY{+w}{ }\PY{o}{=}\PY{+w}{ }\PY{l+s+sa}{raw}\PY{l+s}{\PYZdq{}}\PY{l+s}{Approx }\PY{l+s}{\PYZdq{}}\PY{+w}{ }\PY{o}{*}\PY{+w}{ }\PY{l+s+sa}{raw}\PY{l+s}{\PYZdq{}}\PY{l+s}{\PYZdl{}n = }\PY{l+s}{\PYZdq{}}\PY{+w}{ }\PY{o}{*}\PY{+w}{ }\PY{l+s}{\PYZdq{}}\PY{l+s}{\PYZob{}}\PY{l+s+si}{\PYZdl{}}\PY{p}{(}\PY{n}{n}\PY{p}{)}\PY{l+s}{\PYZcb{}}\PY{l+s}{\PYZdq{}}\PY{+w}{ }\PY{o}{*}\PY{+w}{ }\PY{l+s+sa}{raw}\PY{l+s}{\PYZdq{}}\PY{l+s}{\PYZdl{}}\PY{l+s}{\PYZdq{}}\PY{p}{)}
\PY{n}{scatter!}\PY{p}{(}\PY{n}{ts}\PY{p}{,}\PY{+w}{ }\PY{n}{ys}\PY{p}{,}\PY{+w}{ }\PY{n}{label}\PY{+w}{ }\PY{o}{=}\PY{+w}{ }\PY{l+s+sa}{raw}\PY{l+s}{\PYZdq{}}\PY{l+s}{Data }\PY{l+s}{\PYZdq{}}\PY{+w}{ }\PY{o}{*}\PY{+w}{ }\PY{l+s+sa}{raw}\PY{l+s}{\PYZdq{}}\PY{l+s}{\PYZdl{}n = }\PY{l+s}{\PYZdq{}}\PY{+w}{ }\PY{o}{*}\PY{+w}{ }\PY{l+s}{\PYZdq{}}\PY{l+s}{\PYZob{}}\PY{l+s+si}{\PYZdl{}}\PY{p}{(}\PY{n}{n}\PY{p}{)}\PY{l+s}{\PYZcb{}}\PY{l+s}{\PYZdq{}}\PY{+w}{ }\PY{o}{*}\PY{+w}{ }\PY{l+s+sa}{raw}\PY{l+s}{\PYZdq{}}\PY{l+s}{\PYZdl{}}\PY{l+s}{\PYZdq{}}\PY{p}{,}\PY{+w}{ }\PY{n}{opacity}\PY{+w}{ }\PY{o}{=}\PY{+w}{ }\PY{l+m+mf}{0.5}\PY{p}{)}
\PY{n}{current}\PY{p}{(}\PY{p}{)}
\end{Verbatim}
\end{tcolorbox}

    \begin{center}
    \adjustimage{max size={0.9\linewidth}{0.9\paperheight}}{HW-5_files/HW-5_47_0.pdf}
    \end{center}
    { \hspace*{\fill} \\}
    
    \begin{tcolorbox}[breakable, size=fbox, boxrule=1pt, pad at break*=1mm,colback=cellbackground, colframe=cellborder]
\prompt{In}{incolor}{30}{\boxspacing}
\begin{Verbatim}[commandchars=\\\{\}]
\PY{n}{n}\PY{+w}{ }\PY{o}{=}\PY{+w}{ }\PY{l+m+mi}{80}\PY{p}{;}\PY{+w}{ }
\PY{n}{N}\PY{+w}{ }\PY{o}{=}\PY{+w}{ }\PY{l+m+mi}{2}\PY{+w}{ }\PY{o}{*}\PY{+w}{ }\PY{n}{n}\PY{+w}{ }\PY{o}{+}\PY{+w}{ }\PY{l+m+mi}{1}\PY{p}{;}
\PY{n}{ts}\PY{+w}{ }\PY{o}{=}\PY{+w}{ }\PY{l+m+mi}{2}\PY{o}{*}\PY{p}{(}\PY{o}{\PYZhy{}}\PY{n}{n}\PY{o}{:}\PY{n}{n}\PY{p}{)}\PY{o}{/}\PY{n}{N}
\PY{n}{ys}\PY{+w}{ }\PY{o}{=}\PY{+w}{ }\PY{n}{f}\PY{o}{.}\PY{p}{(}\PY{n}{ts}\PY{p}{)}
\PY{n}{fun}\PY{+w}{ }\PY{o}{=}\PY{+w}{ }\PY{n}{triginterp}\PY{p}{(}\PY{n}{ts}\PY{p}{,}\PY{+w}{ }\PY{n}{ys}\PY{p}{)}
\PY{n}{plot!}\PY{p}{(}\PY{n}{fun}\PY{p}{,}\PY{+w}{ }\PY{o}{\PYZhy{}}\PY{l+m+mi}{1}\PY{p}{,}\PY{+w}{ }\PY{l+m+mi}{1}\PY{p}{,}\PY{+w}{ }\PY{n}{label}\PY{+w}{ }\PY{o}{=}\PY{+w}{ }\PY{l+s+sa}{raw}\PY{l+s}{\PYZdq{}}\PY{l+s}{Approx }\PY{l+s}{\PYZdq{}}\PY{+w}{ }\PY{o}{*}\PY{+w}{ }\PY{l+s+sa}{raw}\PY{l+s}{\PYZdq{}}\PY{l+s}{\PYZdl{}n = }\PY{l+s}{\PYZdq{}}\PY{+w}{ }\PY{o}{*}\PY{+w}{ }\PY{l+s}{\PYZdq{}}\PY{l+s}{\PYZob{}}\PY{l+s+si}{\PYZdl{}}\PY{p}{(}\PY{n}{n}\PY{p}{)}\PY{l+s}{\PYZcb{}}\PY{l+s}{\PYZdq{}}\PY{+w}{ }\PY{o}{*}\PY{+w}{ }\PY{l+s+sa}{raw}\PY{l+s}{\PYZdq{}}\PY{l+s}{\PYZdl{}}\PY{l+s}{\PYZdq{}}\PY{p}{)}
\PY{n}{scatter!}\PY{p}{(}\PY{n}{ts}\PY{p}{,}\PY{+w}{ }\PY{n}{ys}\PY{p}{,}\PY{+w}{ }\PY{n}{mark}\PY{+w}{ }\PY{o}{=}\PY{+w}{ }\PY{l+s+ss}{:o}\PY{p}{,}\PY{+w}{ }\PY{n}{label}\PY{+w}{ }\PY{o}{=}\PY{+w}{ }\PY{l+s+sa}{raw}\PY{l+s}{\PYZdq{}}\PY{l+s}{Data }\PY{l+s}{\PYZdq{}}\PY{+w}{ }\PY{o}{*}\PY{+w}{ }\PY{l+s+sa}{raw}\PY{l+s}{\PYZdq{}}\PY{l+s}{\PYZdl{}n = }\PY{l+s}{\PYZdq{}}\PY{+w}{ }\PY{o}{*}\PY{+w}{ }\PY{l+s}{\PYZdq{}}\PY{l+s}{\PYZob{}}\PY{l+s+si}{\PYZdl{}}\PY{p}{(}\PY{n}{n}\PY{p}{)}\PY{l+s}{\PYZcb{}}\PY{l+s}{\PYZdq{}}\PY{+w}{ }\PY{o}{*}\PY{+w}{ }\PY{l+s+sa}{raw}\PY{l+s}{\PYZdq{}}\PY{l+s}{\PYZdl{}}\PY{l+s}{\PYZdq{}}\PY{p}{,}\PY{+w}{ }\PY{n}{opacity}\PY{+w}{ }\PY{o}{=}\PY{+w}{ }\PY{l+m+mf}{0.5}\PY{p}{)}

\PY{n}{n}\PY{+w}{ }\PY{o}{=}\PY{+w}{ }\PY{l+m+mi}{180}\PY{p}{;}\PY{+w}{ }
\PY{n}{N}\PY{+w}{ }\PY{o}{=}\PY{+w}{ }\PY{l+m+mi}{2}\PY{+w}{ }\PY{o}{*}\PY{+w}{ }\PY{n}{n}\PY{+w}{ }\PY{o}{+}\PY{+w}{ }\PY{l+m+mi}{1}\PY{p}{;}
\PY{n}{ts}\PY{+w}{ }\PY{o}{=}\PY{+w}{ }\PY{l+m+mi}{2}\PY{o}{*}\PY{p}{(}\PY{o}{\PYZhy{}}\PY{n}{n}\PY{o}{:}\PY{n}{n}\PY{p}{)}\PY{o}{/}\PY{n}{N}
\PY{n}{ys}\PY{+w}{ }\PY{o}{=}\PY{+w}{ }\PY{n}{f}\PY{o}{.}\PY{p}{(}\PY{n}{ts}\PY{p}{)}
\PY{n}{fun}\PY{+w}{ }\PY{o}{=}\PY{+w}{ }\PY{n}{triginterp}\PY{p}{(}\PY{n}{ts}\PY{p}{,}\PY{+w}{ }\PY{n}{ys}\PY{p}{)}
\PY{n}{plot!}\PY{p}{(}\PY{n}{fun}\PY{p}{,}\PY{+w}{ }\PY{o}{\PYZhy{}}\PY{l+m+mi}{1}\PY{p}{,}\PY{+w}{ }\PY{l+m+mi}{1}\PY{p}{,}\PY{+w}{ }\PY{n}{label}\PY{+w}{ }\PY{o}{=}\PY{+w}{ }\PY{l+s+sa}{raw}\PY{l+s}{\PYZdq{}}\PY{l+s}{Approx }\PY{l+s}{\PYZdq{}}\PY{+w}{ }\PY{o}{*}\PY{+w}{ }\PY{l+s+sa}{raw}\PY{l+s}{\PYZdq{}}\PY{l+s}{\PYZdl{}n = }\PY{l+s}{\PYZdq{}}\PY{+w}{ }\PY{o}{*}\PY{+w}{ }\PY{l+s}{\PYZdq{}}\PY{l+s}{\PYZob{}}\PY{l+s+si}{\PYZdl{}}\PY{p}{(}\PY{n}{n}\PY{p}{)}\PY{l+s}{\PYZcb{}}\PY{l+s}{\PYZdq{}}\PY{+w}{ }\PY{o}{*}\PY{+w}{ }\PY{l+s+sa}{raw}\PY{l+s}{\PYZdq{}}\PY{l+s}{\PYZdl{}}\PY{l+s}{\PYZdq{}}\PY{p}{)}
\PY{n}{scatter!}\PY{p}{(}\PY{n}{ts}\PY{p}{,}\PY{+w}{ }\PY{n}{ys}\PY{p}{,}\PY{+w}{ }\PY{n}{mark}\PY{+w}{ }\PY{o}{=}\PY{+w}{ }\PY{l+s+ss}{:x}\PY{p}{,}\PY{+w}{ }\PY{n}{label}\PY{+w}{ }\PY{o}{=}\PY{+w}{ }\PY{l+s+sa}{raw}\PY{l+s}{\PYZdq{}}\PY{l+s}{Data }\PY{l+s}{\PYZdq{}}\PY{+w}{ }\PY{o}{*}\PY{+w}{ }\PY{l+s+sa}{raw}\PY{l+s}{\PYZdq{}}\PY{l+s}{\PYZdl{}n = }\PY{l+s}{\PYZdq{}}\PY{+w}{ }\PY{o}{*}\PY{+w}{ }\PY{l+s}{\PYZdq{}}\PY{l+s}{\PYZob{}}\PY{l+s+si}{\PYZdl{}}\PY{p}{(}\PY{n}{n}\PY{p}{)}\PY{l+s}{\PYZcb{}}\PY{l+s}{\PYZdq{}}\PY{+w}{ }\PY{o}{*}\PY{+w}{ }\PY{l+s+sa}{raw}\PY{l+s}{\PYZdq{}}\PY{l+s}{\PYZdl{}}\PY{l+s}{\PYZdq{}}\PY{p}{,}\PY{+w}{ }\PY{n}{opacity}\PY{+w}{ }\PY{o}{=}\PY{+w}{ }\PY{l+m+mf}{0.5}\PY{p}{)}
\end{Verbatim}
\end{tcolorbox}

    \begin{center}
    \adjustimage{max size={0.9\linewidth}{0.9\paperheight}}{HW-5_files/HW-5_48_0.pdf}
    \end{center}
    { \hspace*{\fill} \\}
    
    \hypertarget{q-9.7.3d}{%
\subsection{Q 9.7.3d}\label{q-9.7.3d}}

For integration on a semi-infinite interval such as \([0,\infty)\),
another double exponential transformation is useful:
\(x(t) = \exp(\sinh(t))\).

Write a function \texttt{intsemi(f,tol)} by adapting Function 9.7.3 in
the book for the semi-infinite integration problem. Test it on the
integral

\[\int_0^\infty \frac{e^{-x}}{\sqrt{x}} = \sqrt{\pi}.\]

Compare with the result using \texttt{quadgk} from the QuadGK package.

    You'll need \texttt{intadapt} from Chapter 5, so I've copied it here.

    \begin{tcolorbox}[breakable, size=fbox, boxrule=1pt, pad at break*=1mm,colback=cellbackground, colframe=cellborder]
\prompt{In}{incolor}{31}{\boxspacing}
\begin{Verbatim}[commandchars=\\\{\}]
\PY{k}{function}\PY{+w}{ }\PY{n}{intadapt}\PY{p}{(}\PY{n}{f}\PY{p}{,}\PY{n}{a}\PY{p}{,}\PY{n}{b}\PY{p}{,}\PY{n}{tol}\PY{p}{,}\PY{n}{fa}\PY{o}{=}\PY{n}{f}\PY{p}{(}\PY{n}{a}\PY{p}{)}\PY{p}{,}\PY{n}{fb}\PY{o}{=}\PY{n}{f}\PY{p}{(}\PY{n}{b}\PY{p}{)}\PY{p}{,}\PY{n}{m}\PY{o}{=}\PY{p}{(}\PY{n}{a}\PY{o}{+}\PY{n}{b}\PY{p}{)}\PY{o}{/}\PY{l+m+mi}{2}\PY{p}{,}\PY{n}{fm}\PY{o}{=}\PY{n}{f}\PY{p}{(}\PY{n}{m}\PY{p}{)}\PY{p}{)}
\PY{+w}{    }\PY{c}{\PYZsh{} Use error estimation and recursive bisection.}
\PY{+w}{    }\PY{c}{\PYZsh{} These are the two new nodes and their f\PYZhy{}values.}
\PY{+w}{    }\PY{n}{xl}\PY{+w}{ }\PY{o}{=}\PY{+w}{ }\PY{p}{(}\PY{n}{a}\PY{o}{+}\PY{n}{m}\PY{p}{)}\PY{o}{/}\PY{l+m+mi}{2}\PY{p}{;}\PY{+w}{  }\PY{n}{fl}\PY{+w}{ }\PY{o}{=}\PY{+w}{ }\PY{n}{f}\PY{p}{(}\PY{n}{xl}\PY{p}{)}\PY{p}{;}
\PY{+w}{    }\PY{n}{xr}\PY{+w}{ }\PY{o}{=}\PY{+w}{ }\PY{p}{(}\PY{n}{m}\PY{o}{+}\PY{n}{b}\PY{p}{)}\PY{o}{/}\PY{l+m+mi}{2}\PY{p}{;}\PY{+w}{  }\PY{n}{fr}\PY{+w}{ }\PY{o}{=}\PY{+w}{ }\PY{n}{f}\PY{p}{(}\PY{n}{xr}\PY{p}{)}\PY{p}{;}
\PY{+w}{    }
\PY{+w}{    }\PY{c}{\PYZsh{} Compute the trapezoid values iteratively.}
\PY{+w}{    }\PY{n}{h}\PY{+w}{ }\PY{o}{=}\PY{+w}{ }\PY{p}{(}\PY{n}{b}\PY{o}{\PYZhy{}}\PY{n}{a}\PY{p}{)}
\PY{+w}{    }\PY{n}{T}\PY{+w}{ }\PY{o}{=}\PY{+w}{ }\PY{p}{[}\PY{l+m+mf}{0.}\PY{p}{,}\PY{l+m+mf}{0.}\PY{p}{,}\PY{l+m+mf}{0.}\PY{p}{]}
\PY{+w}{    }\PY{n}{T}\PY{p}{[}\PY{l+m+mi}{1}\PY{p}{]}\PY{+w}{ }\PY{o}{=}\PY{+w}{ }\PY{n}{h}\PY{o}{*}\PY{p}{(}\PY{n}{fa}\PY{o}{+}\PY{n}{fb}\PY{p}{)}\PY{o}{/}\PY{l+m+mi}{2}
\PY{+w}{    }\PY{n}{T}\PY{p}{[}\PY{l+m+mi}{2}\PY{p}{]}\PY{+w}{ }\PY{o}{=}\PY{+w}{ }\PY{n}{T}\PY{p}{[}\PY{l+m+mi}{1}\PY{p}{]}\PY{o}{/}\PY{l+m+mi}{2}\PY{+w}{ }\PY{o}{+}\PY{+w}{ }\PY{p}{(}\PY{n}{h}\PY{o}{/}\PY{l+m+mi}{2}\PY{p}{)}\PY{o}{*}\PY{n}{fm}
\PY{+w}{    }\PY{n}{T}\PY{p}{[}\PY{l+m+mi}{3}\PY{p}{]}\PY{+w}{ }\PY{o}{=}\PY{+w}{ }\PY{n}{T}\PY{p}{[}\PY{l+m+mi}{2}\PY{p}{]}\PY{o}{/}\PY{l+m+mi}{2}\PY{+w}{ }\PY{o}{+}\PY{+w}{ }\PY{p}{(}\PY{n}{h}\PY{o}{/}\PY{l+m+mi}{4}\PY{p}{)}\PY{o}{*}\PY{p}{(}\PY{n}{fl}\PY{o}{+}\PY{n}{fr}\PY{p}{)}
\PY{+w}{    }
\PY{+w}{    }\PY{n}{S}\PY{+w}{ }\PY{o}{=}\PY{+w}{ }\PY{p}{(}\PY{l+m+mi}{4}\PY{n}{T}\PY{p}{[}\PY{l+m+mi}{2}\PY{o}{:}\PY{l+m+mi}{3}\PY{p}{]}\PY{o}{\PYZhy{}}\PY{n}{T}\PY{p}{[}\PY{l+m+mi}{1}\PY{o}{:}\PY{l+m+mi}{2}\PY{p}{]}\PY{p}{)}\PY{+w}{ }\PY{o}{/}\PY{+w}{ }\PY{l+m+mi}{3}\PY{+w}{      }\PY{c}{\PYZsh{} Simpson values}
\PY{+w}{    }\PY{n}{E}\PY{+w}{ }\PY{o}{=}\PY{+w}{ }\PY{p}{(}\PY{n}{S}\PY{p}{[}\PY{l+m+mi}{2}\PY{p}{]}\PY{o}{\PYZhy{}}\PY{n}{S}\PY{p}{[}\PY{l+m+mi}{1}\PY{p}{]}\PY{p}{)}\PY{+w}{ }\PY{o}{/}\PY{+w}{ }\PY{l+m+mi}{15}\PY{+w}{           }\PY{c}{\PYZsh{} error estimate}
\PY{+w}{    }
\PY{+w}{    }\PY{k}{if}\PY{+w}{ }\PY{n}{abs}\PY{p}{(}\PY{n}{E}\PY{p}{)}\PY{+w}{ }\PY{o}{\PYZlt{}}\PY{+w}{ }\PY{n}{tol}\PY{o}{*}\PY{p}{(}\PY{l+m+mi}{1}\PY{o}{+}\PY{n}{abs}\PY{p}{(}\PY{n}{S}\PY{p}{[}\PY{l+m+mi}{2}\PY{p}{]}\PY{p}{)}\PY{p}{)}\PY{+w}{  }\PY{c}{\PYZsh{} acceptable error?}
\PY{+w}{        }\PY{n}{Q}\PY{+w}{ }\PY{o}{=}\PY{+w}{ }\PY{n}{S}\PY{p}{[}\PY{l+m+mi}{2}\PY{p}{]}\PY{+w}{                   }\PY{c}{\PYZsh{} yes\PYZhy{}\PYZhy{}done}
\PY{+w}{        }\PY{n}{nodes}\PY{+w}{ }\PY{o}{=}\PY{+w}{ }\PY{p}{[}\PY{n}{a}\PY{p}{,}\PY{n}{xl}\PY{p}{,}\PY{n}{m}\PY{p}{,}\PY{n}{xr}\PY{p}{,}\PY{n}{b}\PY{p}{]}\PY{+w}{      }\PY{c}{\PYZsh{} all nodes at this level}
\PY{+w}{    }\PY{k}{else}
\PY{+w}{        }\PY{c}{\PYZsh{} Error is too large\PYZhy{}\PYZhy{}bisect and recurse.}
\PY{+w}{        }\PY{n}{QL}\PY{p}{,}\PY{n}{tL}\PY{+w}{ }\PY{o}{=}\PY{+w}{ }\PY{n}{intadapt}\PY{p}{(}\PY{n}{f}\PY{p}{,}\PY{n}{a}\PY{p}{,}\PY{n}{m}\PY{p}{,}\PY{n}{tol}\PY{p}{,}\PY{n}{fa}\PY{p}{,}\PY{n}{fm}\PY{p}{,}\PY{n}{xl}\PY{p}{,}\PY{n}{fl}\PY{p}{)}
\PY{+w}{        }\PY{n}{QR}\PY{p}{,}\PY{n}{tR}\PY{+w}{ }\PY{o}{=}\PY{+w}{ }\PY{n}{intadapt}\PY{p}{(}\PY{n}{f}\PY{p}{,}\PY{n}{m}\PY{p}{,}\PY{n}{b}\PY{p}{,}\PY{n}{tol}\PY{p}{,}\PY{n}{fm}\PY{p}{,}\PY{n}{fb}\PY{p}{,}\PY{n}{xr}\PY{p}{,}\PY{n}{fr}\PY{p}{)}
\PY{+w}{        }\PY{n}{Q}\PY{+w}{ }\PY{o}{=}\PY{+w}{ }\PY{n}{QL}\PY{+w}{ }\PY{o}{+}\PY{+w}{ }\PY{n}{QR}
\PY{+w}{        }\PY{n}{nodes}\PY{+w}{ }\PY{o}{=}\PY{+w}{ }\PY{p}{[}\PY{n}{tL}\PY{p}{;}\PY{n}{tR}\PY{p}{[}\PY{l+m+mi}{2}\PY{o}{:}\PY{k}{end}\PY{p}{]}\PY{p}{]}\PY{+w}{   }\PY{c}{\PYZsh{} merge the nodes w/o duplicate}
\PY{+w}{    }\PY{k}{end}
\PY{+w}{    }\PY{k}{return}\PY{+w}{ }\PY{n}{Q}\PY{p}{,}\PY{n}{nodes}
\PY{k}{end}
\end{Verbatim}
\end{tcolorbox}

    
    \begin{Verbatim}[commandchars=\\\{\}]
intadapt (generic function with 5 methods)
    \end{Verbatim}

    
    \emph{Solution}:

    \begin{tcolorbox}[breakable, size=fbox, boxrule=1pt, pad at break*=1mm,colback=cellbackground, colframe=cellborder]
\prompt{In}{incolor}{32}{\boxspacing}
\begin{Verbatim}[commandchars=\\\{\}]
\PY{k}{function}\PY{+w}{ }\PY{n}{intsemi}\PY{p}{(}\PY{n}{f}\PY{p}{,}\PY{n}{tol}\PY{p}{)}\PY{+w}{   }
\PY{+w}{    }\PY{n}{x}\PY{+w}{ }\PY{o}{=}\PY{+w}{ }\PY{n}{t}\PY{+w}{ }\PY{o}{\PYZhy{}\PYZgt{}}\PY{+w}{ }\PY{n}{exp}\PY{p}{(}\PY{n}{sinh}\PY{p}{(}\PY{n}{t}\PY{p}{)}\PY{p}{)}
\PY{+w}{    }\PY{n}{dx\PYZus{}dt}\PY{+w}{ }\PY{o}{=}\PY{+w}{ }\PY{n}{t}\PY{+w}{ }\PY{o}{\PYZhy{}\PYZgt{}}\PY{+w}{ }\PY{n}{cosh}\PY{p}{(}\PY{n}{t}\PY{p}{)}\PY{o}{*}\PY{n}{exp}\PY{p}{(}\PY{n}{sinh}\PY{p}{(}\PY{n}{t}\PY{p}{)}\PY{p}{)}
\PY{+w}{    }\PY{n}{g}\PY{+w}{ }\PY{o}{=}\PY{+w}{ }\PY{n}{t}\PY{+w}{ }\PY{o}{\PYZhy{}\PYZgt{}}\PY{+w}{ }\PY{n}{f}\PY{p}{(}\PY{n}{x}\PY{p}{(}\PY{n}{t}\PY{p}{)}\PY{p}{)}\PY{o}{*}\PY{n}{dx\PYZus{}dt}\PY{p}{(}\PY{n}{t}\PY{p}{)}

\PY{+w}{    }\PY{c}{\PYZsh{} Find where to truncate the integration interval.}
\PY{+w}{    }\PY{n}{M}\PY{+w}{ }\PY{o}{=}\PY{+w}{ }\PY{l+m+mi}{3}
\PY{+w}{    }\PY{k}{while}\PY{+w}{ }\PY{p}{(}\PY{n}{abs}\PY{p}{(}\PY{n}{g}\PY{p}{(}\PY{o}{\PYZhy{}}\PY{n}{M}\PY{p}{)}\PY{p}{)}\PY{+w}{ }\PY{o}{\PYZgt{}}\PY{+w}{ }\PY{n}{tol}\PY{o}{/}\PY{l+m+mi}{100}\PY{p}{)}\PY{+w}{ }\PY{o}{||}\PY{+w}{ }\PY{p}{(}\PY{n}{abs}\PY{p}{(}\PY{n}{g}\PY{p}{(}\PY{n}{M}\PY{p}{)}\PY{p}{)}\PY{+w}{ }\PY{o}{\PYZgt{}}\PY{+w}{ }\PY{n}{tol}\PY{o}{/}\PY{l+m+mi}{100}\PY{p}{)}
\PY{+w}{        }\PY{n}{M}\PY{+w}{ }\PY{o}{+=}\PY{+w}{ }\PY{l+m+mf}{0.5}
\PY{+w}{        }\PY{k}{if}\PY{+w}{ }\PY{n}{isinf}\PY{p}{(}\PY{n}{x}\PY{p}{(}\PY{n}{M}\PY{p}{)}\PY{p}{)}\PY{+w}{ }
\PY{+w}{            }\PY{n+nd}{@warn}\PY{+w}{ }\PY{l+s}{\PYZdq{}}\PY{l+s}{Function may not decay fast enough.}\PY{l+s}{\PYZdq{}}
\PY{+w}{            }\PY{n}{M}\PY{+w}{ }\PY{o}{\PYZhy{}=}\PY{+w}{ }\PY{l+m+mf}{0.5}
\PY{+w}{            }\PY{k}{break}
\PY{+w}{        }\PY{k}{end}
\PY{+w}{    }\PY{k}{end}

\PY{+w}{    }\PY{n}{I}\PY{p}{,}\PY{n}{t}\PY{+w}{ }\PY{o}{=}\PY{+w}{ }\PY{n}{intadapt}\PY{p}{(}\PY{n}{g}\PY{p}{,}\PY{o}{\PYZhy{}}\PY{n}{M}\PY{p}{,}\PY{n}{M}\PY{p}{,}\PY{n}{tol}\PY{p}{)}
\PY{+w}{	}\PY{k}{return}\PY{+w}{ }\PY{n}{I}\PY{p}{,}\PY{n}{x}\PY{o}{.}\PY{p}{(}\PY{n}{t}\PY{p}{)}
\PY{k}{end}
\end{Verbatim}
\end{tcolorbox}

    
    \begin{Verbatim}[commandchars=\\\{\}]
intsemi (generic function with 1 method)
    \end{Verbatim}

    
    \begin{tcolorbox}[breakable, size=fbox, boxrule=1pt, pad at break*=1mm,colback=cellbackground, colframe=cellborder]
\prompt{In}{incolor}{33}{\boxspacing}
\begin{Verbatim}[commandchars=\\\{\}]
\PY{n}{res}\PY{+w}{ }\PY{o}{=}\PY{+w}{ }\PY{n}{intsemi}\PY{p}{(}\PY{p}{(}\PY{n}{x}\PY{p}{)}\PY{+w}{ }\PY{o}{\PYZhy{}\PYZgt{}}\PY{+w}{ }\PY{n}{exp}\PY{p}{(}\PY{o}{\PYZhy{}}\PY{n}{x}\PY{p}{)}\PY{+w}{ }\PY{o}{/}\PY{+w}{ }\PY{n}{sqrt}\PY{p}{(}\PY{n}{x}\PY{p}{)}\PY{p}{,}\PY{+w}{ }\PY{l+m+mf}{1e\PYZhy{}6}\PY{p}{)}
\end{Verbatim}
\end{tcolorbox}

    
    \begin{Verbatim}[commandchars=\\\{\}]
(1.772455041233968, [2.853912066288294e-20, 1.047305076987121e-17, 1.7716071483390152e-15, 1.529189428251244e-13, 7.357070993478746e-12, 4.1997533690540406e-11, 2.130199545856566e-10, 9.677925382088462e-10, 3.96790228465493e-9, 1.4783835383662816e-8  …  11.24834714080576, 13.605631155782989, 20.712703825724148, 33.48319690621568, 58.02575654373648, 108.9804279373336, 4006.68761532299, 2.180430367273218e6, 1.3592365778261388e11, 3.5039621991597236e19])
    \end{Verbatim}

    
    \begin{tcolorbox}[breakable, size=fbox, boxrule=1pt, pad at break*=1mm,colback=cellbackground, colframe=cellborder]
\prompt{In}{incolor}{34}{\boxspacing}
\begin{Verbatim}[commandchars=\\\{\}]
\PY{n}{sqrt}\PY{p}{(}\PY{n+nb}{pi}\PY{p}{)}
\end{Verbatim}
\end{tcolorbox}

    
    \begin{Verbatim}[commandchars=\\\{\}]
1.7724538509055159
    \end{Verbatim}

    
    \begin{tcolorbox}[breakable, size=fbox, boxrule=1pt, pad at break*=1mm,colback=cellbackground, colframe=cellborder]
\prompt{In}{incolor}{35}{\boxspacing}
\begin{Verbatim}[commandchars=\\\{\}]
\PY{c}{\PYZsh{} verify}
\PY{n}{res}\PY{p}{[}\PY{l+m+mi}{1}\PY{p}{]}\PY{+w}{ }\PY{o}{\PYZhy{}}\PY{+w}{ }\PY{n}{sqrt}\PY{p}{(}\PY{n+nb}{pi}\PY{p}{)}
\end{Verbatim}
\end{tcolorbox}

    
    \begin{Verbatim}[commandchars=\\\{\}]
1.1903284520276003e-6
    \end{Verbatim}

    

    % Add a bibliography block to the postdoc
    
    
    
\end{document}
